% ===========================================================================
%  AVL Trees --- CSE 200
%  A 9-minute, three-part talk. Beamer + TikZ only; every figure is computed
%  and drawn natively (no external plotting language).
%
%  Compile:
%    pdflatex -interaction=nonstopmode avl.tex   (twice, for the frame count)
%
%  Numbers on these slides are simulated ground truth from verify/ --- see
%  verify/avl_verify.ps1 and verify/cascade_sweep.ps1. Do not edit a figure
%  here without re-running those.
%
%  Design: avltheme.sty implements the deck's design spec (1280x720 px canvas,
%  1 px = 0.0125 cm = 0.3543 pt). Read a value from the spec, convert, place.
% ===========================================================================
\documentclass[aspectratio=169,10pt,xcolor={table}]{beamer}

\usepackage[T1]{fontenc}
\usepackage{amsmath,amssymb}
\usepackage{tikz}
\usepackage{pgfplots}
\pgfplotsset{compat=1.18}
\usetikzlibrary{arrows.meta,positioning,calc,decorations.pathreplacing,
                shapes.geometric,fit,backgrounds}

% Load order matters: fibtree ships safe standalone defaults, the theme
% overrides them for the light ground. Theme must come second.
% ---------------------------------------------------------------------------
%  Recursive minimum-node AVL ("Fibonacci") tree macro.
%
%  T_h is drawn from T_{h-1} and T_{h-2} by the SAME recursion the proof uses,
%  and the horizontal layout is computed from phi via Binet's formula - so the
%  golden ratio literally positions the picture that proves the golden ratio.
%
%  Leaf count of T_h is L(h) = F(h+1); subtree centres are offset by half the
%  sibling's leaf count, which packs the tree with no overlap at any height.
%
%  Usage inside a tikzpicture:
%      \setavlscale{<leaf sep>}{<level sep>}
%      \drawfib{7}{0}{0}
% ---------------------------------------------------------------------------

\newcounter{avlnode}
\newcommand{\avlleafsep}{0.55}
\newcommand{\avllevsep}{0.80}

\newcommand{\setavlscale}[2]{%
  \renewcommand{\avlleafsep}{#1}%
  \renewcommand{\avllevsep}{#2}%
}

% Leaves of T_h via Binet: L(h) = F(h+1). Returns 0 for h < 0 (empty subtree).
\newcommand{\avlL}[1]{round((((1+sqrt(5))/2)^((#1)+1))/sqrt(5))}

% Node styles (overridden by the theme; these are safe standalone defaults)
\tikzset{
  avldot/.style   = {circle, draw=black, fill=black, inner sep=0pt, minimum size=3pt},
  avlgolddot/.style={circle, draw=black, fill=black, inner sep=0pt, minimum size=3pt},
  avledge/.style  = {draw=black, line width=0.4pt},
}

% \drawfib{height}{x}{y}
% Wrapped in a group so each frame's \avlself survives its own recursive calls;
% the counter is global, so node names stay unique across the whole picture.
\newcommand{\drawfib}[3]{%
  \begingroup
  \stepcounter{avlnode}%
  \edef\avlself{f\theavlnode}%
  \node[avldot] (\avlself) at (#2,#3) {};%
  \ifnum#1>0\relax
    % ---- left child is T(h-1); offset by half the RIGHT sibling's leaves ----
    \pgfmathsetmacro{\avlRsib}{\avlL{#1-2}}%
    \pgfmathsetmacro{\avlcxA}{(#2)-(\avlRsib/2)*\avlleafsep}%
    \pgfmathsetmacro{\avlcyA}{(#3)-\avllevsep}%
    \edef\avlnexth{\the\numexpr#1-1\relax}%
    \edef\avlchildA{f\the\numexpr\value{avlnode}+1\relax}%
    % fully expand the coordinates INTO the call, otherwise the recursive
    % invocation would re-bind \avlcxA before it is ever used
    \edef\avlgo{\noexpand\drawfib{\avlnexth}{\avlcxA}{\avlcyA}}\avlgo
    \draw[avledge] (\avlself) -- (\avlchildA);%
  \fi
  \ifnum#1>1\relax
    % ---- right child is T(h-2); offset by half the LEFT sibling's leaves ----
    \pgfmathsetmacro{\avlLsib}{\avlL{#1-1}}%
    \pgfmathsetmacro{\avlcxB}{(#2)+(\avlLsib/2)*\avlleafsep}%
    \pgfmathsetmacro{\avlcyB}{(#3)-\avllevsep}%
    \edef\avlnexthb{\the\numexpr#1-2\relax}%
    \edef\avlchildB{f\the\numexpr\value{avlnode}+1\relax}%
    \edef\avlgob{\noexpand\drawfib{\avlnexthb}{\avlcxB}{\avlcyB}}\avlgob
    \draw[avledge] (\avlself) -- (\avlchildB);%
  \fi
  \endgroup
}

% Convenience: reset the id counter before each tree so names are predictable
\newcommand{\resetavl}{\setcounter{avlnode}{0}}

% ---------------------------------------------------------------------------
%  Motion cues for rotation frames.
%
%  Beamer overlays are jump-cuts by nature -- one frame replaces the last.
%  These macros fake continuous motion across that cut with two tricks:
%    1. GHOSTING: redraw the pivot's PRE-rotation position at low opacity
%       behind the post-rotation tree, connected by a faint dashed path.
%       Flipping between the two overlay steps then reads as the ghost
%       "solidifying" into its new spot, instead of a hard swap.
%    2. PIVOT ARCS: a curved arrow swept around the pivot node showing the
%       rotation's direction, so which way "up" happens is never ambiguous.
% ---------------------------------------------------------------------------

% \rotarc{pivot coord}{radius}{start angle}{end angle}
% Curved arrow around a pivot, drawn in the Fibonacci-gold rotation colour.
\newcommand{\rotarc}[4]{%
  \draw[avlgoldarrow, line width=1.2pt]
    (#1) ++(#3:#2) arc (#3:#4:#2);%
}
% Same, in the alert colour -- for a deletion-triggered rotation.
\newcommand{\rotarcalert}[4]{%
  \draw[draw=avlalert, -{Stealth[length=4pt]}, line width=1.2pt]
    (#1) ++(#3:#2) arc (#3:#4:#2);%
}

% \pivotdot{coord} -- small solid marker pinning down the exact pivot point.
\newcommand{\pivotdot}[1]{%
  \fill[avlfib] (#1) circle (1.6pt);%
}

% \ghostnode{x}{y}{label} -- a faded afterimage of a node at its OLD spot.
\newcommand{\ghostnode}[3]{%
  \node[circle, draw=avlmuted, fill=avlground, text=avlmuted, opacity=0.32,
        inner sep=1pt, minimum size=6.5mm, font=\footnotesize] at (#1,#2) {#3};%
}
% \ghostpath{from}{to} -- faint dashed thread linking an old spot to the new one.
\newcommand{\ghostpath}[2]{%
  \draw[avlmuted, opacity=0.35, dashed, line width=0.5pt] (#1) -- (#2);%
}

\usepackage{avltheme}

% Allow building one part alone for solo rehearsal:
%   \includeonlyframes{p1}   % or p2 / p3
%\includeonlyframes{p1}

\begin{document}

% ===========================================================================
%  PART 1 --- Introduction, the BST problem, the AVL idea, height, proof,
%             Fibonacci property.
%  Follows the requested flow:
%    intro -> what is a BST -> what problem -> AVL idea -> how AVL solves it
%    -> height property -> proof -> build-up -> golden ratio -> exact count
%  Wording rule: short sentences, common words, no idioms.
%
%  Height convention throughout the deck: h counts EDGES, so a single node is
%  height 0 and N(0)=1. A lecture that calls a single node "height 1" gets
%  N(h) = F(h+2)-1; ours is F(h+3)-1. Same fact, shifted index. Do not mix them.
% ===========================================================================

% ------------------------------------------------------------------ title
\begin{frame}[plain,label=p1]
  % 5 px three-hue top bar: the three parts, in the order they arrive.
  % Anchored to the page itself, not to the text block, so it bleeds edge to
  % edge; the same picture carries the progress mark at the opposite corner.
  \begin{tikzpicture}[remember picture,overlay]
    \fill[avlpone]   (current page.north west) rectangle
                     ($(current page.north west)+(0.3333\paperwidth,-5\pxl)$);
    \fill[avlptwo]   ($(current page.north west)+(0.3333\paperwidth,0)$) rectangle
                     ($(current page.north west)+(0.6666\paperwidth,-5\pxl)$);
    \fill[avlpthree] ($(current page.north west)+(0.6666\paperwidth,0)$) rectangle
                     ($(current page.north east)+(0,-5\pxl)$);
    \fill[avlpone]   (current page.south west) rectangle
                     ($(current page.south west)+
                       (\dimexpr\paperwidth/\inserttotalframenumber\relax,3\pxl)$);
    % Course mark sits on the footer line the content frames use: same 84 px
    % gutter, clear of the progress tick in the corner below it.
    \node[anchor=west,inner sep=0pt]
      at ($(current page.south west)+(84\pxl,42\pxl)$) {\eyebrow{CSE 200}};
  \end{tikzpicture}
  \vfill
  \begin{center}
    {\avlfDeck\wxbold\color{avlink}AVL Trees}\\[18\pxl]
    {\avlfLead\wmed\color{avlfib}Trees which balance all on their own}\\[34\pxl]
    % Fibonacci glyph, flanked by hairlines: the thread, planted before a word
    % of it is spoken. leafSep 13 px, levSep 16 px, per the spec.
    \begin{tikzpicture}[x=1\pxl,y=1\pxl,
        avldot/.style={circle,draw=avlfib!62!avlground,
                       fill=avlfib!62!avlground,inner sep=0pt,minimum size=6.4\pxl},
        avledge/.style={draw=avlfib!42!avlground,line width=1.4\pxl}]
      \setavlscale{13}{16}
      \resetavl
      \drawfib{7}{0}{-46}
      \draw[avlrule,line width=1\pxl] (-243,-46) -- (-155,-46);
      \draw[avlrule,line width=1\pxl] ( 155,-46) -- ( 243,-46);
    \end{tikzpicture}\\[30\pxl]
    % Three equal-width columns, so the roll centres under its own name and the
    % blocks space evenly however long the names are. The roll takes the
    % eyebrow treatment -- faint and tracked -- so it reads as a subtitle to the
    % name rather than competing with it.
    \begin{minipage}[b]{0.30\textwidth}\centering
      {\avlfBody\wmed\color{avlbody}Md Ishtiak Moin}\\[7\pxl]
      {\avlfEyebrow\wsemi\color{avlfaint}\avltrack{2305011}}
    \end{minipage}%
    \begin{minipage}[b]{0.30\textwidth}\centering
      {\avlfBody\wmed\color{avlbody}Sheikh Raufur Rahim}\\[7\pxl]
      {\avlfEyebrow\wsemi\color{avlfaint}\avltrack{2305010}}
    \end{minipage}%
    \begin{minipage}[b]{0.30\textwidth}\centering
      {\avlfBody\wmed\color{avlbody}Abhijnan Podder}\\[7\pxl]
      {\avlfEyebrow\wsemi\color{avlfaint}\avltrack{2305007}}
    \end{minipage}
  \end{center}
  \vfill
\end{frame}

\avldivider{1}{The problem}{When can a search tree stop being fast?}{avlpone}{avlponeA}{avlponeB}{avlponeC}

\setavlpart{1}{Part 1 --- The problem}{avlpone}

% ------------------------------------------------------------------ what is a BST
\begin{frame}[label=p1]{First, the rule every BST follows}
  \vspace{0.3em}
  \begin{columns}[T,onlytextwidth]
    \begin{column}{0.52\textwidth}
      {\avlfLead\wbold\color{avlink}Binary Search Tree}\\[0.7em]
      {\color{avlbody}For \emph{every} node:}
      \begin{itemize}
        \setlength{\itemsep}{0.35em}
        \item everything on the \fib{left} is \fib{smaller}
        \item everything on the \fib{right} is \fib{bigger}
      \end{itemize}
      \vspace{0.6em}
      {\avlfSec\color{avlmuted}To search, you compare once and drop one whole
       side.}\\[0.9em]
      \begin{avlpanel*}
        Cost of a search $=$ \fib{the height of the tree}.
      \end{avlpanel*}
    \end{column}
    \begin{column}{0.44\textwidth}
      \begin{center}
      \begin{tikzpicture}[x=1cm,y=1cm,scale=0.98]
        \node[avlkey] (n4) at (0,0)       {4};
        \node[avlkey] (n2) at (-1.35,-1.2) {2};
        \node[avlkey] (n6) at (1.35,-1.2)  {6};
        \node[avlkey] (n1) at (-2.1,-2.4)  {1};
        \node[avlkey] (n3) at (-0.6,-2.4)  {3};
        \node[avlkey] (n5) at (0.6,-2.4)   {5};
        \node[avlkey] (n7) at (2.1,-2.4)   {7};
        \foreach \a/\b in {n4/n2,n4/n6,n2/n1,n2/n3,n6/n5,n6/n7}
          \draw[avledge] (\a) -- (\b);
        \node[avlannot] at (-2.05,-0.50) {smaller};
        \node[avlannot] at ( 2.05,-0.50) {bigger};
        % annotation row under the figure: hairline, key, hairline
        \draw[avlrule,line width=1\pxl] (-2.55,-3.15) -- (-0.95,-3.15);
        \draw[avlrule,line width=1\pxl] ( 0.95,-3.15) -- ( 2.55,-3.15);
        \node[color=avlfib,font=\avlfSec\wmed] at (0,-3.15) {height $=2$};
      \end{tikzpicture}
      \end{center}
    \end{column}
  \end{columns}

  \note{Say the rule out loud. Left is smaller, right is bigger. That is all a
  BST is. And because you throw away half the tree at each step, the cost of a
  search is just the height.}
\end{frame}

% ------------------------------------------------------------------ the problem
\begin{frame}[label=p1]{The problem: sorted input destroys it}
  % Captions sit ABOVE each tree, a hairline divides the two halves, and the
  % colour carries the verdict: green is the legal-and-short tree, red the list.
  % The green nodes on the left are the search path to 7 -- three of them.
  \vspace{0.1em}
  \begin{center}
  \begin{tikzpicture}[x=1cm,y=1cm,scale=0.94,every node/.style={transform shape}]
    \figdivider{0}{0.95}{-3.75}

    \begin{scope}[xshift=-3.55cm,yshift=-0.55cm]
      \node[avlannot,text=avlmuted] at (0,1.35) {inserted 4, 2, 6, 1, 3, 5, 7};
      \node[avlkeyok]  (b4) at (0,0)        {4};
      \node[avlkey]    (b2) at (-1.25,-1.0) {2};
      \node[avlkeyok]  (b6) at (1.25,-1.0)  {6};
      \node[avlkey]    (b1) at (-1.95,-2.0) {1};
      \node[avlkey]    (b3) at (-0.55,-2.0) {3};
      \node[avlkey]    (b5) at (0.55,-2.0)  {5};
      \node[avlkeyok]  (b7) at (1.95,-2.0)  {7};
      \foreach \a/\b in {b4/b2,b2/b1,b2/b3,b6/b5}
        \draw[avledge] (\a) -- (\b);
      \foreach \a/\b in {b4/b6,b6/b7}
        \draw[draw=avlbalance,line width=1.6\pxl] (\a) -- (\b);
      % both verdicts share one baseline: this scope sits 0.40 lower than the
      % right-hand one, so its label y is 0.40 higher to compensate
      \node[color=avlbalance,font=\avlfStmt\wbold] at (0,-2.95) {3 steps};
    \end{scope}

    \begin{scope}[xshift=3.55cm,yshift=-0.15cm]
      \node[avlannot,text=avlmuted] at (0,0.95) {inserted 1, 2, 3, 4, 5, 6, 7};
      % step must exceed the disc diameter or the edges vanish inside the nodes
      \foreach \k [count=\i from 0] in {1,2,3,4,5,6,7} {
        \pgfmathsetmacro{\yy}{-\i*0.44}
        \pgfmathsetmacro{\xx}{-1.20+\i*0.40}
        \node[avlkeyalert,minimum size=34\pxl,font=\avlfAnnot] (c\k) at (\xx,\yy) {\k};
      }
      \foreach \a/\b in {c1/c2,c2/c3,c3/c4,c4/c5,c5/c6,c6/c7}
        \draw[draw=avlalert,line width=1.6\pxl] (\a) -- (\b);
      \node[color=avlalert,font=\avlfStmt\wbold] at (0,-3.35) {7 steps};
    \end{scope}
  \end{tikzpicture}
  \end{center}

  \vspace{0.4em}
  \begin{center}
    {\avlfSec\color{avlmuted}Same seven keys. Both are legal BSTs.
     The right one is now a \viol{list}, not a tree.}
  \end{center}

  \note{Both trees are legal. The right one is what you get when the data
  arrives already sorted --- IDs, dates, imported records. This is not a rare
  case. It is the most common case. Search cost goes from log n to n.}
\end{frame}

% ------------------------------------------------------------------ the idea
\begin{frame}[label=p1]{The idea: check the height at every node}
  \vspace{0.6em}
  \begin{center}
    {\avlfLead\color{avlbody}The tree got tall because one side grew and
     \emph{nobody checked}.}
  \end{center}
  \vspace{0.6em}
  \begin{center}
  \begin{tikzpicture}[x=1cm,y=1cm,scale=0.90,every node/.style={transform shape}]
    \node[avlkey,minimum size=42\pxl] (r) at (0,0) {};
    % subtrees as ember triangles, each carrying its height as the symbol
    \node[avltri,minimum size=2.3cm,font=\avlfLead,text=avlfib] (L) at (-1.75,-1.55) {$a$};
    \node[avltri,minimum size=1.75cm,font=\avlfLead,text=avlfib] (R) at (1.75,-1.25) {$b$};
    \draw[avledge] (r) -- (L);
    \draw[avledge] (r) -- (R);
    % height rulers, one per side, measured from the root row to each base
    \heightbrace{-3.05}{-0.30}{-2.32}{height $a$}{avlfaint}
    \heightbrace{3.05}{-0.30}{-1.75}{height $b$}{avlfaint}
    \draw[avlrule,line width=1\pxl,dotted] (-3.05,-0.30) -- (3.05,-0.30);
    \draw[avlrule,line width=1\pxl] (-2.55,-2.72) -- (2.55,-2.72);
    \node[avlannot,anchor=north] at (0,-2.80) {one node, two subtree heights};
  \end{tikzpicture}
  \end{center}

  \vspace{0.5em}
  \begin{center}
    \begin{avlpanel}[0.62\textwidth]
      The rule: $a$ and $b$ may differ by at most \fib{\textbf{1}}
    \end{avlpanel}
  \end{center}

  \note{The fix is simple to state. At every node, look at the height of the
  left side and the height of the right side. Do not let them differ by more
  than one. That is the whole idea.}
\end{frame}

% ------------------------------------------------------------------ DEFINITION
\begin{frame}[label=p1]{Definition: what an AVL tree is}
  \vspace{0.2em}
  \begin{columns}[T,onlytextwidth]
    \begin{column}{0.55\textwidth}
      \vspace{0.1em}
      {\avlfSec\color{avlmuted}Adelson-Velsky and Landis, 1962. The first
       self-balancing search tree.}\\[0.8em]

      \begin{avlpanel*}
        Balance factor of a node\\[6\pxl]
        $\mathrm{bf} = (\text{height of left}) - (\text{height of right})$
      \end{avlpanel*}\\[0.7em]

      {\strong{An AVL tree is:}}
      \begin{itemize}
        \setlength{\itemsep}{0.25em}
        \item a binary search tree, \emph{and}
        \item every node has \ok{$\mathrm{bf} = -1$, $0$, or $+1$}
      \end{itemize}
      \vspace{0.4em}
      {\avlfSec\color{avlmuted}If any node reaches $+2$ or $-2$, the tree is
       \viol{broken} and must be fixed.}
    \end{column}
    \begin{column}{0.42\textwidth}
      \begin{center}
      \begin{tikzpicture}[x=1cm,y=1cm,scale=0.92,every node/.style={transform shape}]
        \node[avlkeyok] (a) at (0,0) {};
        \node[avlannotok,anchor=west] at (0.40,0.16) {$+1$};
        \node[avlkey] (b) at (-1.15,-1.25) {};
        \node[avlannotok,anchor=east] at (-1.52,-1.06) {$0$};
        \node[avlkey] (c) at (1.15,-1.25) {};
        \node[avlannotok,anchor=west] at (1.52,-1.06) {$0$};
        \node[avlkey] (d) at (-1.9,-2.5) {};
        \node[avlkey] (e) at (-0.4,-2.5) {};
        \foreach \p/\q in {a/b,a/c,b/d,b/e} \draw[avledge] (\p) -- (\q);
        \node[avlannot] at (0,-3.30) {left height 2, right height 1};
        \node[color=avlbalance,font=\avlfSec] at (0,-3.85) {legal AVL tree};
      \end{tikzpicture}
      \end{center}
    \end{column}
  \end{columns}

  \note{Now the name and the rule. AVL stands for the two Soviet
  mathematicians who invented it in 1962. Balance factor is left height minus
  right height. An AVL tree is a BST where every single node has a balance
  factor of minus one, zero, or plus one.}
\end{frame}

% ------------------------------------------------------------------ how it solves it
\begin{frame}[label=p1]{How AVL fixes the bad case}
  \vspace{-0.2em}
  \begin{center}
  \begin{tikzpicture}[x=1cm,y=1cm,scale=0.86,every node/.style={transform shape}]
    \begin{scope}[xshift=-3.6cm,yshift=0.4cm]
      \foreach \k [count=\i from 0] in {1,2,3,4,5,6,7} {
        \pgfmathsetmacro{\yy}{-\i*0.55}
        \pgfmathsetmacro{\xx}{\i*0.30}
        \node[avlkeyalert,minimum size=36\pxl,font=\avlfAnnot] (c\k) at (\xx,\yy) {\k};
      }
      \foreach \a/\b in {c1/c2,c2/c3,c3/c4,c4/c5,c5/c6,c6/c7}
        \draw[draw=avlalert,line width=1.6\pxl] (\a) -- (\b);
      \node[avlannot,text=avlmuted] at (0.9,-4.2) {plain BST};
      \node[color=avlalert,font=\avlfLead\wbold] at (0.9,-4.85) {height 6};
    \end{scope}

    % neutral arrow: the transformation is not part of the Fibonacci thread
    \draw[avlarrow] (-0.9,-1.6) -- (0.5,-1.6);
    \node[avlannot,above=2pt] at (-0.2,-1.55) {same input};

    \begin{scope}[xshift=3.4cm,yshift=-0.6cm]
      \node[avlkeyok] (a4) at (0,0)        {4};
      \node[avlkeyok] (a2) at (-1.3,-1.15) {2};
      \node[avlkeyok] (a6) at (1.3,-1.15)  {6};
      \node[avlkeyok] (a1) at (-2.0,-2.3)  {1};
      \node[avlkeyok] (a3) at (-0.6,-2.3)  {3};
      \node[avlkeyok] (a5) at (0.6,-2.3)   {5};
      \node[avlkeyok] (a7) at (2.0,-2.3)   {7};
      \foreach \a/\b in {a4/a2,a4/a6,a2/a1,a2/a3,a6/a5,a6/a7}
        \draw[draw=avlbalance!70,line width=1.6\pxl] (\a) -- (\b);
      \node[avlannot,text=avlmuted] at (0,-3.2) {AVL tree};
      \node[color=avlbalance,font=\avlfLead\wbold] at (0,-3.85) {height 2};
    \end{scope}
  \end{tikzpicture}
  \end{center}

  \vspace{0.7em}
  \begin{center}
    {\avlfSec\color{avlmuted}Insert 1, 2, 3, 4, 5, 6, 7 into both.
     AVL rebalances as it goes. \ok{7 steps become 3.}}
  \end{center}
  \vspace{0.5em}

  \note{Same sorted input that destroyed the plain tree. The AVL tree fixes
  itself while you insert, and stays short. We will see exactly how it does
  that in Part 2.}
\end{frame}

% ------------------------------------------------------------------ height property
\begin{frame}[label=p1]{The height property}
  \vspace{0.1em}
  \begin{center}
    {\avlfLead\color{avlbody}A local rule at every node gives a
     \emph{global} guarantee:}\\[0.9em]
    \begin{avlpanel}[0.50\textwidth]
      {\avlfDivider\color{avlfib}$h < 1.44 \, \log_2 n$}
    \end{avlpanel}
  \end{center}

  % Three columns, a single hairline under the header, no rules between rows
  % and no fills: the payoff cell is carried by colour alone.
  \vspace{0.9em}
  \begin{center}
  \setlength{\tabcolsep}{20\pxl}
  \renewcommand{\arraystretch}{1.45}
  \setlength{\arrayrulewidth}{1\pxl}
  \arrayrulecolor{avlrule}
  {\avlfSec
  \begin{tabular}{lcc}
    \color{avlfaint}nodes & \color{avlfaint}perfect tree
      & \color{avlfaint}worst AVL tree\\\hline
    $1\,000\,000$ & 20 steps & \fib{\textbf{28 steps}}\\
  \end{tabular}}
  \end{center}

  \vspace{0.7em}
  \begin{center}
    {\avlfSec\color{avlmuted}Never a list. Always close to the best possible.}
  \end{center}

  \note{This is the promise. No matter what order the data arrives in, the
  height stays within 44 percent of a perfect tree. One million records: at
  most 28 comparisons. Next question --- why 1.44? Where does that number come
  from?}
\end{frame}

% ------------------------------------------------------------------ proof setup
%  The framing beat only. The recurrence itself is NOT stated here any more --
%  it is built from actual small trees on the next slide, so it arrives as
%  something the audience has already watched happen four times.
\begin{frame}[label=p1]{Proof: ask the question backwards}
  \vspace{0.9em}
  \begin{center}
    {\avlfLead\color{avlmuted}Instead of ``how tall can it get?'', ask:}\\[0.8em]
    % the pivot line: the turned question, with the turn itself in ink
    {\avlfStmt\color{avlfib}\wbold What is the
     {\color{avlink}\mdseries smallest} AVL tree of height $h$?}
  \end{center}

  \vspace{0.7em}
  \centerline{{\color{avlrule}\rule[2\pxl]{70\pxl}{1\pxl}}\hskip10\pxl
              {\color{avlfaint}\rule{4\pxl}{4\pxl}}\hskip10\pxl
              {\color{avlrule}\rule[2\pxl]{70\pxl}{1\pxl}}}

  \vspace{0.7em}
  \begin{center}
    {\avlfBody\color{avlbody}If even the emptiest legal tree of height $h$
     needs a lot of nodes,\\[0.3em]
     then a tree with few nodes \emph{cannot} be tall.}\\[0.9em]
    \begin{avlpanel}[0.58\textwidth]
      Write \fib{$N(h)$} for that smallest node count.
    \end{avlpanel}
  \end{center}

  \note{Turn the question around. Instead of bounding the tallest tree, build
  the emptiest legal tree of a given height and count it. If the emptiest tree
  of height h already needs a lot of nodes, then a tree with few nodes cannot
  possibly be that tall. Same fact, and this direction we can actually count.}
\end{frame}

% ------------------------------------------------------------------ BUILD-UP
%  The four smallest minimal trees, revealed one at a time. Each new tree is
%  visibly the two before it, joined under one root -- so the recurrence is
%  SEEN before it is written down.
%
%  Coordinate map (per panel, origin at the panel's root):
%    T_0  a(0,0)
%    T_1  a(0,0)  b(-0.45,-0.8)
%    T_2  a(0,0)  b(-0.62,-0.8) c(-1.05,-1.6)   d(0.62,-0.8)
%    T_3  a(0,0)  b(-0.85,-0.8) c(-1.35,-1.6) d(-1.75,-2.4) e(-0.42,-1.6)
%                 f(0.85,-0.8)  g(0.45,-1.6)
%  Panel origins: -5.0, -2.9, 0.0, 3.8. Captions on a common baseline.
%  NOTE: the style key must not be called "cap" -- that collides with TikZ's
%  own line-cap key and the frame will not compile.
% Title kept to a single line on purpose: a second line pushes the rule down
% and steals the height the ladder and its panel need.
\begin{frame}[label=p1]{Each tree is the two before it, under one root}
  \vspace{-0.7em}
  \begin{center}
  % Vertical axis only (y=0.86cm), and no `transform shape`: the ladder loses
  % height so the panel below gets air, while the horizontal pitch stays at
  % 1 cm -- the T_k / N(h) captions are set at full size and would collide with
  % each other if the columns closed up. Node discs shrink by their own key.
  \begin{tikzpicture}[x=1cm,y=0.86cm,
      mini/.style   = {circle, draw=avlfib, fill=avlfib!12!avlground,
                       line width=1.3\pxl, inner sep=0pt, minimum size=4.0mm},
      miniroot/.style={circle, draw=avlfib, fill=avlfib, inner sep=0pt,
                       line width=1.3\pxl, minimum size=4.0mm},
      medge/.style  = {avlminiedge},
      blk/.style    = {draw=avlfaint, dashed, line width=1\pxl,
                       rounded corners=2pt, inner sep=2.2pt},
      blab/.style   = {avlannot},
      treecap/.style= {font=\avlfSec, text=avlink, align=center,
                       minimum width=2cm, minimum height=8mm}]
    % fixed extent so nothing shifts between overlays
    % Extent trimmed to what the beats actually occupy (deepest caption sits at
    % -3.85), so the panel below gets its room without shrinking the figure.
    \path (-5.9,-3.95) rectangle (5.5,0.30);

    % ---------------- T_0 ----------------
    \begin{scope}[xshift=-5.0cm]
      \node[miniroot] (za) at (0,0) {};
      \node[treecap] at (0,-3.25) {$T_0$\\ \fib{$N(0)=1$}};
    \end{scope}

    % ---------------- T_1 ----------------
    \only<2->{\begin{scope}[xshift=-2.9cm]
      \node[miniroot] (oa) at (0,0) {};
      \node[mini] (ob) at (-0.45,-0.8) {};
      \draw[medge] (oa) -- (ob);
      \node[treecap] at (0,-3.25) {$T_1$\\ \fib{$N(1)=2$}};
    \end{scope}}

    % ---------------- T_2 = root + T_1 + T_0 ----------------
    \only<3->{\begin{scope}[xshift=0cm]
      \node[miniroot] (ta) at (0,0) {};
      \node[mini] (tb) at (-0.62,-0.8) {};
      \node[mini] (tc) at (-1.05,-1.6) {};
      \node[mini] (td) at ( 0.62,-0.8) {};
      \draw[medge] (ta) -- (tb);  \draw[medge] (tb) -- (tc);
      \draw[medge] (ta) -- (td);
      \node[blk, fit=(tb)(tc), label={[blab]below:$T_1$}] {};
      \node[blk, fit=(td),     label={[blab]below:$T_0$}] {};
      \node[treecap] at (0,-3.25) {$T_2$\\ \fib{$N(2)=1+2+1=4$}};
    \end{scope}}

    % ---------------- T_3 = root + T_2 + T_1 ----------------
    \only<4->{\begin{scope}[xshift=3.8cm]
      \node[miniroot] (ha) at (0,0) {};
      \node[mini] (hb) at (-0.85,-0.8) {};
      \node[mini] (hc) at (-1.35,-1.6) {};
      \node[mini] (hd) at (-1.75,-2.4) {};
      \node[mini] (he) at (-0.42,-1.6) {};
      \node[mini] (hf) at ( 0.85,-0.8) {};
      \node[mini] (hg) at ( 0.45,-1.6) {};
      \draw[medge] (ha) -- (hb);  \draw[medge] (hb) -- (hc);
      \draw[medge] (hc) -- (hd);  \draw[medge] (hb) -- (he);
      \draw[medge] (ha) -- (hf);  \draw[medge] (hf) -- (hg);
      \node[blk, fit=(hb)(hc)(hd)(he), label={[blab]below:$T_2$}] {};
      \node[blk, fit=(hf)(hg),         label={[blab]below:$T_1$}] {};
      \node[treecap] at (0,-3.45) {$T_3$\\ \fib{$N(3)=1+4+2=7$}};
    \end{scope}}
  \end{tikzpicture}
  \end{center}

  % The caption block reserves its own height, so no beat shifts the figure.
  \begin{center}
  \begin{minipage}[t][112\pxl][c]{0.86\textwidth}\centering
    \only<1>{{\avlfSec\color{avlmuted}The smallest legal tree of height 0 is
              one node.}}
    \only<2>{{\avlfSec\color{avlmuted}For height 1 you need a child --- and one
              child is enough, because a gap of 1 is legal.}}
    \only<3>{{\avlfSec\color{avlmuted}For height 2, one side must reach
              height 1. The other is as \emph{small} as the rule allows:
              height 0.}}
    \only<4>{{\avlfSec\color{avlmuted}And again: one side $T_2$, the other side
              $T_1$. \fib{The pattern does not change.}}}
    % the recurrence arrives in a panel -- it is the payoff, not a caption
    \only<5->{\begin{avlpanel}[0.74\textwidth]
                {\avlfTitle\color{avlfib}$N(h) = N(h-1) + N(h-2) + 1$}\\[8\pxl]
                {\avlfEyebrow\color{avlmuted}one root, plus the two smallest
                 trees that fit under it}
              \end{avlpanel}}
  \end{minipage}
  \end{center}

  \note{Click through them slowly. Height zero: one node. Height one: add a
  child. Height two: the left side has to reach height one, and the right side
  is as small as the rule permits, which is height zero. Point at the two
  dashed boxes. Height three: same again --- a T two and a T one. By now they
  should be saying the recurrence before you write it. Then click and write it.}
\end{frame}

% ------------------------------------------------------- GOLDEN RATIO (algebra)
\begin{frame}[label=p1]{Where the golden ratio comes from}
  % Beats 1-4 build the algebra top-down: a faint eyebrow announces each move,
  % the maths follows in serif, and the line that matters turns ember. Beat 5
  % clears the board for the payoff -- phi itself.
  \only<1-4>{
  \vspace{-0.8em}
  \begin{center}
    {\avlfSec\color{avlmuted}The $+1$ is in the way --- so absorb it.}\\[0.6em]
    {\avlfLead $N(h) = N(h-1) + N(h-2) + 1$}\\[0.6em]

    \onslide<2->{%
      \eyebrow{add 1 to both sides, and write $M(h) = N(h)+1$}\\[0.4em]
      {\avlfBody $\big(N(h)+1\big) = \big(N(h{-}1)+1\big) + \big(N(h{-}2)+1\big)$}}
    \\[0.7em]

    \onslide<3->{%
      {\avlfLead\color{avlfib}$M(h) = M(h-1) + M(h-2)$}\\[0.25em]
      {\avlfEyebrow\color{avlmuted}Fibonacci exactly, no $+1$ left}}
  \end{center}

  \onslide<4->{
  \vspace{0.1em}
  \begin{center}
    {\color{avlrule}\rule{0.72\textwidth}{1\pxl}}\\[0.3em]
    \eyebrow{look for solutions of the form $M(h) = x^h$}\\[0.25em]
    {\avlfBody $x^h = x^{h-1} + x^{h-2}
      \;\Longrightarrow\; x^2 = x + 1
      \;\Longrightarrow\; x = (1 \pm \sqrt5)\,/\,2$}
  \end{center}}}

  % Beat 5 -- the phi payoff. A split panel: the glyph, a rule, the value.
  \only<5>{
  \vspace{0.1em}
  \begin{center}
    \eyebrow{$x^2 = x+1$ has one root above 1, and it dominates}\\[1.0em]
    % Laid out in tikz so the glyph, the rule and the two value lines all
    % share one vertical centre -- baseline stacking put phi low in the box.
    % The glyph and its value, filling the panel: phi set large and bold in
    % ink, the number in ember so the two are told apart at a glance.
    \begin{avlpanel}[0.38\textwidth]
      \begin{tikzpicture}[baseline=(mid.base)]
        \coordinate (mid) at (0,-0.08);
        \node[anchor=east,inner sep=0pt] at (-0.30,0)
          {\fontsize{50}{50}\selectfont\color{avlink}$\boldsymbol{\varphi}$};
        \draw[avlrule,line width=1\pxl] (0,-0.46) -- (0,0.46);
        \node[anchor=west,inner sep=0pt] at (0.30,0)
          {\fontsize{26}{26}\selectfont\wbold\color{avlfib}1.618};
      \end{tikzpicture}
    \end{avlpanel}\\[0.9em]
    {\avlfTitle\color{avlink}$N(h) \sim \varphi^{\,h}$}\\[0.9em]
    \begin{minipage}{0.74\textwidth}\centering
      {\avlfBody\color{avlbody}The smallest legal tree multiplies by
       \fib{1.618 every level}. Nodes grow exponentially, so height can only
       grow like a logarithm --- \fib{this is the whole reason 1.44 appears in
       this talk.}}
    \end{minipage}
  \end{center}}

  \note{The recurrence is almost Fibonacci but for the plus one. Adding one to
  both sides makes it vanish --- M of h is exactly Fibonacci. Then the standard
  trick: guess a geometric solution, divide through by x to the h minus two,
  and you are left with x squared equals x plus one. That is the defining
  equation of the golden ratio. So the smallest AVL tree grows like phi to the
  h --- and that is the only reason 1.44 appears anywhere in this talk.}
\end{frame}

% --------------------------------------------------------- the exact identity
\begin{frame}[label=p1]{The exact count}
  % The identity is set bare, big and ember -- it needs no box; the panel at
  % the foot of the frame is where the consequence lands.
  \vspace{0.2em}
  \begin{center}
    {\avlfDivider\color{avlfib}$N(h) = F(h+3) - 1$}\\[0.7em]
    {\avlfSec\color{avlmuted}because $M(h) = N(h)+1$ is Fibonacci, starting at
     $M(0) = 2 = F(3)$}
  \end{center}

  % No vertical rules, no zebra: a hairline under the header, avlrule2 between
  % body rows, and the payoff column carries the ember tint.
  \vspace{0.7em}
  \onslide<2->{
  \begin{center}
  \setlength{\tabcolsep}{10\pxl}
  \renewcommand{\arraystretch}{1.35}
  \setlength{\arrayrulewidth}{1\pxl}
  {\avlfSec
  \begin{tabular}{r|cccccccc}
    \color{avlfaint}$h$ & 0 & 1 & 2 & 3 & 4 & 5 & 6 & \fib{\textbf{7}}\\
    \color{avlfaint}$N(h)$ & 1 & 2 & 4 & 7 & 12 & 20 & 33 & \fib{\textbf{54}}\\
    \color{avlfaint}$F(h{+}3)$ & 2 & 3 & 5 & 8 & 13 & 21 & 34 & \fib{\textbf{55}}\\
  \end{tabular}}
  \end{center}}

  \vspace{0.8em}
  \onslide<3->{
  \begin{center}
    \begin{avlpanel}[0.58\textwidth]
      \panellabel{taking logs of $N(h) \sim \varphi^h$}
      {\avlfTitle\color{avlfib}$h \approx \log_2 n \,/\, \log_2\varphi
        = 1.44 \log_2 n$}
    \end{avlpanel}
  \end{center}}

  \note{One line. The plus-one substitution we just made IS the minus one here.
  Read the last column out loud: at height seven the smallest AVL tree has 54
  nodes, and 55 is a Fibonacci number. Hold onto 54 --- Part 3 will build this
  exact tree and break it.}
\end{frame}

% ------------------------------------------------------------------ THE PLANT
\begin{frame}[label=p1]{Remember this number: 54}
  \vspace{-0.4em}
  \begin{center}
  % height-7 Fibonacci tree: leafSep 44 px, levSep 38 px (spec section 5)
  \begin{tikzpicture}[x=1cm,y=1cm,
      avldot/.style={circle,draw=avlfib,fill=avlfib,inner sep=0pt,
                     minimum size=5.8\pxl},
      avledge/.style={draw=avlfib!75,line width=1.3\pxl}]
    \setavlscale{0.48}{0.37}
    \resetavl
    \drawfib{7}{0}{0}
  \end{tikzpicture}
  \end{center}

  \vspace{0.1em}
  \begin{center}
    {\avlfLead\color{avlink}This is the \fib{worst tree AVL allows} at
     height 7.}\\[0.4em]
    {\avlfSec\color{avlmuted}It has the fewest nodes possible:
     $N(7) = F(10) - 1 = 55 - 1 =$ \fib{\textbf{54}}}\\[0.7em]
    % the Part 3 plant, in Part 3's hue -- the only plum on a Part 1 slide
    \begin{avlpanelh}[0.98\textwidth]{avlpthree}
      {\color{avlpthree}In Part 3 we will delete one leaf from this exact tree
       --- and watch it come apart.}
    \end{avlpanelh}
  \end{center}

  \note{This is the emptiest legal AVL tree of height seven. Fibonacci predicts
  it has 54 nodes. Remember that number. In Part 3 we will take one leaf out of
  it and count the damage. Handoff: we know the tree stays short. Now, how does
  it fix itself?}
\end{frame}

% ===========================================================================
%  PART 2 --- Rotations and Insertion
%  Flow: one rotation drawn as a physical turn -> the four cases table ->
%        the 2x2 grid of pivots -> the insertion story (which ends on a
%        question) -> one simulated insert that answers it -> the catch ->
%        why one rotation is always enough, and what it costs
%
%  Rotation language, used identically on every rotation in this deck:
%    a PIN goes into the node the rotation is NAMED at (the broken one), and
%    one big ARC over the tree shows which way the whole thing turns --
%    right rotation = clockwise, left rotation = anticlockwise.
%  The result is always a SEPARATE overlay, with ghosts at the old positions.
% ===========================================================================

\avldivider{2}{Rotations}{How the tree fixes itself.}{avlptwo}{avlptwoA}{avlptwoB}{avlptwoC}

\setavlpart{2}{Part 2 --- Rotations}{avlptwo}

% ------------------------------------------------------------------ ONE ROTATION
%  The whole idea in one frame: pin a node, turn the tree about it.
%
%  Layout rule for this frame: a node's x is its position in sorted order and
%  NEVER changes; only its depth does. A rotation therefore moves nodes
%  straight up or straight down, which is what makes the ghosts readable --
%  every faded circle sits directly above or below where its node landed. It
%  also makes the payoff visible rather than asserted: nothing moves sideways,
%  so the search order cannot break.
\begin{frame}[label=p2]{A rotation turns the tree about a pin}
  \vspace{-0.2em}
  \begin{center}
  \begin{tikzpicture}[x=1cm,y=1cm,scale=0.95,every node/.style={transform shape}]
    % fixed extent: the picture must not jump between overlay steps. The top is
    % set by the arc (radius 1.65 about a node at y = 0.8).
    \path (-3.9,-2.78) rectangle (3.9,2.55);

    % ---------------- before: 30, 20, 10 all went left ----------------
    \only<1-2>{
      \draw[avledge] (1.9,0.8) -- (0,-0.5) -- (-1.9,-1.8);
      \node[avlkeyalert] at (1.9,0.8)   {30};
      \node[avlkey]      at (0,-0.5)    {20};
      \node[avlkeygold]  at (-1.9,-1.8) {10};
    }
    \only<1>{%
      \draw[avlalert,line width=2\pxl] (1.9,0.8) circle (0.42);
      \node[avlannotalert,anchor=west] at (2.42,0.98) {$\mathrm{bf} = +2$};}
    % Pin and turn are ONE beat: the pin is only meaningful as the thing the
    % arc sweeps about, and splitting them made the audience click twice for
    % half a gesture.
    \only<2>{\pivotpin{1.9,0.8}\rotright{1.9,0.8}{1.65}}

    % ---------------- after: 20 came up, 30 came down -----------------
    \only<3->{
      % Afterimages on the result beat ONLY. They carry the eye across the jump
      % cut from the turn; by the next beat the order strip is making the point
      % instead, and three dashed discs would just read as extra nodes.
      \only<3>{%
        \ghostnode{1.9}{0.8}{30}  \ghostpath{1.9,0.45}{1.9,-0.15}
        \ghostnode{0}{-0.5}{20}   \ghostpath{0,-0.15}{0,0.45}
        \ghostnode{-1.9}{-1.8}{10} \ghostpath{-1.9,-1.45}{-1.9,-0.85}}
      % restored: the tree that comes out of the turn is legal, so it is green
      \draw[draw=avlbalance!70,line width=1.6\pxl] (-1.9,-0.5) -- (0,0.8) -- (1.9,-0.5);
      \node[avlkeyok] at (0,0.8)     {20};
      \node[avlkeyok] at (-1.9,-0.5) {10};
      \node[avlkeyok] at (1.9,-0.5)  {30};
      % No pin on these beats. The pin was a fixed post; redrawing it under 30
      % after the turn would say it travelled with the node, which is the one
      % thing the whole picture is trying not to say.
    }

    % ---------------- the order strip, from the payoff beat on --------
    \only<4->{
      \draw[avlrule,line width=1\pxl] (-2.6,-2.20) -- (2.6,-2.20);
      \foreach \k/\x in {10/-1.9, 20/0, 30/1.9} {
        \draw[avlrule,line width=1\pxl] (\x,-2.08) -- (\x,-2.32);
        \node[color=avlink,font=\avlfAnnot] at (\x,-1.90) {\k};
      }
      \node[avlannot,anchor=north] at (0,-2.42)
        {left to right order --- \strong{never changed}};
    }
  \end{tikzpicture}
  \end{center}

  \vspace{-0.3em}
  \begin{center}
    \only<1>{{\avlfSec\color{avlmuted}Left skewed tree.
              \viol{30 has a balance factor of $+2$}.}}
    \only<2>{{\avlfSec\color{avlmuted}Push the pin into
              {\color{avlptwo}\wmed 30} and turn the whole tree
              {\color{avlptwo}\wmed clockwise} about it. That is a
              {\color{avlptwo}\wmed right rotation}.}}
    \only<3>{{\avlfSec\color{avlmuted}20 became the parent, 30 came down on the
              right. \ok{The tree is balanced again.}}}
    \only<4->{{\avlfSec\color{avlmuted}Nothing moved sideways, so the order is
              safe, and only \fib{three links} changed.}}
  \end{center}

  \note{Say it as a physical action. The pin goes into 30, the node that is
  broken, and in the same breath you turn the whole tree clockwise about that
  pin, the way you would turn a stiff wire. One gesture, not two.
  20 ends up on top, 30 ends up on the right. Point at
  the faded circles: that is where each node was a second ago. Then point at
  the ruler: nothing slid left or right. That is why a rotation can never break
  the search rule, and it only ever re-links three pointers, so it is constant
  time.}
\end{frame}
% ------------------------------------------------------------------ four cases
\begin{frame}[label=p2]{Four cases, and what to do}
  % No vertical rules, no zebra. Header labels avlfaint, hairline under the
  % header, avlrule2 between body rows, the fix column carries the ember tint.
  % One hairline under the header and nothing else: no vertical rules, no row
  % rules, no fills. The fix column is carried by colour alone.
  \vspace{0.1em}
  \begin{center}
  \setlength{\tabcolsep}{20\pxl}
  \renewcommand{\arraystretch}{1.32}
  \setlength{\arrayrulewidth}{1\pxl}
  \arrayrulecolor{avlrule}
  {\avlfBody
  \begin{tabular}{lll}
    {\avlfEyebrow\color{avlfaint}case}
      & {\avlfEyebrow\color{avlfaint}where the new node went}
      & {\avlfEyebrow\color{avlfaint}fix}\\\hline
    \strong{LL} & left child's \emph{left} side  & \fib{one right rotation}\\
    \strong{RR} & right child's \emph{right} side & \fib{one left rotation}\\
    \strong{LR} & left child's \emph{right} side & \fib{left, then right}\\
    \strong{RL} & right child's \emph{left} side & \fib{right, then left}\\
  \end{tabular}}
  \end{center}

  \vspace{0.5em}
  \begin{center}
    {\avlfSec\color{avlmuted}Each case is named after the two steps the new
     node took down from the broken node.}\\[0.7em]
    \begin{avlpanel}[0.96\textwidth]
      A \fib{zig-zag} imbalance needs two rotations: first on the child, then
      on the broken node.
    \end{avlpanel}
  \end{center}

  \note{Four shapes, four fixes. The two straight ones need a single rotation.
  The two bent ones need two rotations: first straighten the bend, then rotate
  as normal. Do not linger --- the next slide shows all four side by side.}
\end{frame}

% ------------------------------------------------------------------ 2x2 GRID
%  The four cases as four actual pivots, laid out 2x2. Same three keys every
%  time (10, 20, 30) so the only variable is arrival ORDER, and every cell ends
%  on the same balanced tree.
%
%  Cell shapes are written once here: the four cases are just the same three
%  nodes in four arrangements, and the doubles turn into the singles after
%  their first rotation, which is the point the frame is making.
% Cell lattice, the same rule as the simulation frame: 10 always sits at
% x = -0.90, 20 at x = 0, 30 at x = +0.90, in every shape. A click therefore
% moves a node straight up or straight down and never sideways, and the bent
% shapes visibly straighten into the chains directly above them.
\newcommand{\cellchainL}{%   leans left -- the LL shape
  \draw[avledge,line width=0.9pt] (0.90,0) -- (0,-0.80) -- (-0.90,-1.60);
  \node[gred] at (0.90,0) {30}; \node[gkey] at (0,-0.80) {20};
  \node[ggold] at (-0.90,-1.60) {10};}
\newcommand{\cellchainR}{%   leans right -- the RR shape
  \draw[avledge,line width=0.9pt] (-0.90,0) -- (0,-0.80) -- (0.90,-1.60);
  \node[gred] at (-0.90,0) {10}; \node[gkey] at (0,-0.80) {20};
  \node[ggold] at (0.90,-1.60) {30};}
% The two bent shapes do NOT use the chain's even level drop. A bend spans two
% lattice columns on its upper edge and only one on its lower edge, so at an
% even 0.80/0.80 drop the upper edge comes out 1.97 long against a 1.20 stub --
% which is what made LR and RL read lopsided next to the chains. The column
% positions are fixed by sorted order and cannot move, so the only free lever
% is the drop: 0.52 then 1.08 keeps the total height at 1.60 (the tags and the
% row hairline never move) and brings the two edges to 1.87 against 1.42.
% The bend node is the ONLY node on its level in these shapes, so an uneven
% drop is invisible -- there is nothing beside it to measure against.
\newcommand{\cellbentL}{%    left then right -- the LR shape
  \draw[avledge,line width=0.9pt] (0.90,0) -- (-0.90,-0.52) -- (0,-1.60);
  \node[gred] at (0.90,0) {30}; \node[gkey] at (-0.90,-0.52) {10};
  \node[ggold] at (0,-1.60) {20};}
\newcommand{\cellbentR}{%    right then left -- the RL shape
  \draw[avledge,line width=0.9pt] (-0.90,0) -- (0.90,-0.52) -- (0,-1.60);
  \node[gred] at (-0.90,0) {10}; \node[gkey] at (0.90,-0.52) {30};
  \node[ggold] at (0,-1.60) {20};}
\newcommand{\cellflat}{%     the balanced result every case lands on
  \draw[avlgoldedge,line width=1.2pt] (-0.90,-0.80) -- (0,0) -- (0.90,-0.80);
  \node[ggold] at (0,0) {20}; \node[ggold] at (-0.90,-0.80) {10};
  \node[ggold] at (0.90,-0.80) {30};}

\begin{frame}[label=p2]{\only<1>{Four cases, four pivots}%
    \only<2>{LR and RL just become LL and RR}%
    \only<3->{Four arrivals, one tree}}
  \vspace{-1.1em}
  \begin{center}
  % The grid used to sit at scale 0.70 inside a 14 cm text block -- barely 8 cm
  % wide, with the discs and their tags shrunk to match. The lattice is now
  % drawn larger and spread wider, so the figure fills the column it is given.
  % The grid used to sit at scale 0.70 inside a 14 cm text block -- barely 8 cm
  % wide, with the discs and their tags shrunk to match. The lattice is drawn
  % larger and spread wider now, and each tag is ONE line: two-line tags were
  % what made the block too tall to enlarge without pushing into the panel.
  \begin{tikzpicture}[x=1cm,y=1cm,scale=0.72,every node/.style={transform shape},
      gkey/.style ={avlkey,      minimum size=6.6mm,font=\avlfEyebrow,inner sep=0pt},
      ggold/.style={avlkeygold,  minimum size=6.6mm,font=\avlfEyebrow,inner sep=0pt},
      gred/.style ={avlkeyalert, minimum size=6.6mm,font=\avlfEyebrow,inner sep=0pt},
      tag/.style  ={font=\avlfBody,align=center,text=avlmuted,anchor=north,
                    inner sep=1pt}]
    % hairline between the two rows of cells, clear of the row-1 tags above it
    % and the row-2 discs below it
    \only<3->{\draw[avlrule,line width=1\pxl] (-6.6,-2.68) -- (6.6,-2.68);}
    \setpinscale{1.0}
    % Fixed extent, symmetric about x = 0. Row pitch 3.32: an arc reaches 0.85
    % above its cell origin and a one-line tag reaches 2.34 below it, so
    % anything tighter puts the top row's caption inside the bottom row's arc.
    \path (-7.1,-5.66) rectangle (7.1,0.88);

    % ================= LL : 30, 20, 10 =================
    \begin{scope}[shift={(-5.2,0)}]
      % arc first, so the nodes always sit on top of it
      \only<1>{\rotright{0.90,0}{0.85}\cellchainL\pivotpin{0.90,0}}
      \only<2->{\cellflat}
      \node[tag] at (0,-2.02) {\textbf{\color{avlink}LL}\ \ 30, 20, 10
        \ \quietc{---}\ \fib{right 30}};
    \end{scope}

    % ================= RR : 10, 20, 30 =================
    \begin{scope}[shift={(5.2,0)}]
      \only<1>{\rotleft{-0.90,0}{0.85}\cellchainR\pivotpin{-0.90,0}}
      \only<2->{\cellflat}
      \node[tag] at (0,-2.02) {\textbf{\color{avlink}RR}\ \ 10, 20, 30
        \ \quietc{---}\ \fib{left 10}};
    \end{scope}

    % ================= LR : 30, 10, 20 =================
    % The first turn is at 10 and straightens the bend into an LL chain; the
    % second turn is then literally the LL cell sitting above it.
    \begin{scope}[shift={(-5.2,-3.32)}]
      \only<1>{\rotleft{-0.90,-0.52}{0.70}
               \cellbentL\pivotpin{-0.90,-0.52}}
      \only<2>{\rotright{0.90,0}{0.85}\cellchainL\pivotpin{0.90,0}}
      \only<3->{\cellflat}
      \node[tag] at (0,-2.02) {\textbf{\color{avlink}LR}\ \ 30, 10, 20
        \ \quietc{---}\ \fib{left, right}};
    \end{scope}

    % ================= RL : 10, 30, 20 =================
    \begin{scope}[shift={(5.2,-3.32)}]
      \only<1>{\rotright{0.90,-0.52}{0.70}
               \cellbentR\pivotpin{0.90,-0.52}}
      \only<2>{\rotleft{-0.90,0}{0.85}\cellchainR\pivotpin{-0.90,0}}
      \only<3->{\cellflat}
      \node[tag] at (0,-2.02) {\textbf{\color{avlink}RL}\ \ 10, 30, 20
        \ \quietc{---}\ \fib{right, left}};
    \end{scope}
  \end{tikzpicture}
  \end{center}

  \only<3->{
  \vspace{0.45em}
  \begin{center}
    \begin{avlpanel}[0.96\textwidth]
      Same three keys, four arrival orders --- \fib{all four land on the same
      tree}.
    \end{avlpanel}
  \end{center}
  \vspace{0.15em}}

  \note{Four cells, same three keys, only the arrival order changes. Point at
  each pin and each arc: LL turns clockwise about 30, RR turns anticlockwise
  about 10. Click. LL and RR are done. Now the important beat: the bent ones
  are no longer bent --- LR has turned into an LL chain and RL into an RR
  chain. Click again. Same pin, same turn as the cell directly above, and all
  four land on the same tree.}
\end{frame}

% ------------------------------------------------------------------ rotation power
\begin{frame}[label=p2]{The beauty: one constant operation fixes everything}
  \vspace{-0.4em}
  \begin{center}
  \eyebrow{What makes a rotation elegant}
  \end{center}

  \vspace{0.5em}
  \begin{center}
  \begin{minipage}{0.86\textwidth}
  \begin{columns}[T,onlytextwidth]
    \begin{column}{0.47\textwidth}
      {\fib{Only three links change}}\\[0.15em]
      {\avlfSec\color{avlmuted}$O(1)$ no matter the tree size}\\[0.8em]
      {\fib{Restores subtree height}}\\[0.15em]
      {\avlfSec\color{avlmuted}Parent sees no change, repair stops}
    \end{column}
    \begin{column}{0.47\textwidth}
      {\fib{Search order is safe}}\\[0.15em]
      {\avlfSec\color{avlmuted}Left-right order never breaks}\\[0.8em]
      {\fib{One rotation suffices}}\\[0.15em]
      {\avlfSec\color{avlmuted}No need to check nodes above}
    \end{column}
  \end{columns}
  \end{minipage}
  \end{center}

  \vspace{1.0em}
  \begin{center}
    \begin{avlpanel}[0.84\textwidth]
      Three pointers, no recursion, no loops --- that is why a rotation
      \fib{runs in constant time}.
    \end{avlpanel}\\[0.7em]
    {\avlfSec\color{avlmuted}Four cases. \quad One pattern. \quad
     \ok{Perfect balance.}}
  \end{center}

  \note{The heart of AVL. A rotation is O(1): only three pointers move. It
  restores the subtree to its original height, so the parent tree sees nothing
  has changed. That is why one rotation is always enough for insertion. This is
  the elegant core that makes AVL practical.}
\end{frame}

% ------------------------------------------------------------------ insertion story
%  The insertion arc, four frames:
%    1. the story        insert -> walk up -> rotate, ending on the QUESTION
%    2. the simulation   one insert on a tree nobody has seen, answering it
%    3. the catch        one rotation? always? yes
%    4. why so           the height argument, then the cost
%
%  The question is asked before it is answered on purpose. The old deck
%  asserted "one rotation is always enough" on the rule slide and then ran a
%  simulation that could only confirm what had already been given away.
%
%  WORDING, fixed across all three frames up to the catch: the third step is
%  "rotate when we find a bad node". Not "the first bad node" -- that phrasing
%  quietly promises there is only one, which is the thing the catch is about
%  to ask.
\begin{frame}[label=p2]{Insertion: insert, walk up, rotate}
  \vspace{0.7em}
  \begin{center}
  \begin{tikzpicture}[x=1cm,y=1cm,
      stepbox/.style={draw=avlptwo!45!avlground,fill=avlptwo!8!avlground,
                   line width=1.5\pxl,rounded corners=6pt,
                   minimum width=3.4cm,minimum height=1.15cm,
                   font=\avlfSec,text=avlink},
      last/.style={draw=avlfib!45!avlground,fill=avlfib!8!avlground,
                   line width=1.5\pxl,rounded corners=6pt,
                   minimum width=3.4cm,minimum height=1.15cm,
                   font=\avlfSec,text=avlink}]
    \node[stepbox] (s1) at (-4.4,0)
      {\begin{minipage}{3.0cm}\centering Insert like a normal BST\end{minipage}};
    \node[stepbox] (s2) at (0,0)
      {\begin{minipage}{3.0cm}\centering Walk back up to the root\end{minipage}};
    \node[last] (s3) at (4.4,0)
      {\begin{minipage}{3.0cm}\centering Rotate when we find a bad node\end{minipage}};
    \draw[avlgoldarrow] (s1) -- (s2);
    \draw[avlgoldarrow] (s2) -- (s3);
  \end{tikzpicture}
  \end{center}

  \vspace{0.9em}
  \begin{center}
  % Fixed-height box so the flow diagram above does not shift when the
  % question replaces the caption.
  \begin{minipage}[t][150\pxl][t]{0.96\textwidth}
  \centering
    \only<1>{{\avlfSec\color{avlmuted}The tree is only repaired on the way
              \emph{back up}.}\\[0.9em]
             {\avlfSec\color{avlmuted}A node is bad once its balance factor
              reaches $\pm 2$.}}
    \only<2->{%
      \begin{avlpanel}[0.78\textwidth]
        {\avlfLead\color{avlfib}\wbold But exactly how many rotations?}
      \end{avlpanel}\\[0.8em]
      {\avlfSec\color{avlmuted}Let us run one insert and count.}}
  \end{minipage}
  \end{center}

  \note{Three steps, and the order matters: you insert like an ordinary BST,
  then you walk back up, and when you find a node with a balance factor of plus
  or minus two, you rotate there. Click. Now hold the question open: how many
  rotations does that cost? Do not answer it yet, the next slide answers it by
  counting.}
\end{frame}
% ------------------------------------------------------------------ INSERTION SIM
%  ONE insert into a tree nobody has seen before, so the count is earned rather
%  than asserted. The tree is chosen so the walk up passes two healthy nodes
%  before it finds the broken one: the whole point of step 2 is that most of the
%  walk finds nothing, and a tree that broke at the parent would hide that.
%
%    start        40(20(10,30),60)          legal: bf(40)=+1, bf(20)=0
%    insert 5     lands left of 10          bf(10)=+1 ok, bf(20)=+1 ok,
%                                           bf(40)=+2 BROKEN, an LL shape
%    fix          one right rotation at 40  -> 20(10(5),40(30,60))
%
%  Height check, which is the reason one rotation ends it: the tree stood at
%  height 2 before the insert and stands at height 2 after the rotation.
%
%  LAYOUT: lattice pitch 1.25 x 1.20 at scale 0.84, the same figures the 1..7
%  simulation used, so this tree reads at the size the deck has already
%  established. A key's COLUMN is its position in sorted order and never
%  changes, only its ROW does. c0..c5 = 5, 10, 20, 30, 40, 60.
%
%  There is no caption strip under this figure. Two things were competing for
%  the same 1 cm: a readable tree and a line of prose. The tree won, and what
%  the prose was carrying (the balance factors, the shape name, the count) is
%  annotated on the figure itself where it belongs.
%
%  ALGORITHM RAIL. All four steps hold their slots from beat one, so the list
%  is a map of where we are rather than a list being built. \setbeamercovered
%  does the dimming: a step outside its own overlays is covered, which beamer
%  renders at 32% opacity, and \alt paints it muted grey underneath. The green
%  chevron is drawn only on the active beats. "Done" is the exception: it is
%  \only-gated, so it is absent, not dimmed, until the tree is actually legal.

% \insmark{<active spec>} -- the "you are here" chevron. Always occupies its
% slot (the \path reserves the box) and is only inked on the active beats, so
% the rail cannot shift sideways between overlays. 1 unit = 1 design px.
\newcommand{\insmark}[1]{%
  \begin{tikzpicture}[baseline=0pt,x=0.0125cm,y=0.0125cm]
    \path (0,0) rectangle (11,12);
    \only<#1>{\fill[avlbalance] (0,1) -- (0,11) -- (9.5,6) -- cycle;}
  \end{tikzpicture}}

% \insstep[<reveal spec>]{<active spec>}{<text>}
\newcommand{\insstep}[3][1-]{%
  \hbox to\linewidth{%
    \insmark{#2}\hskip14\pxl
    \begin{minipage}[t]{\dimexpr\linewidth-40\pxl\relax}
      \avlfSec\raggedright\strut
      \only<#1>{\uncover<#2>{%
        \alt<#2>{\color{avlbalance}\wsemi}{\color{avlmuted}}#3}}%
    \end{minipage}\hfil}%
  \vskip 24\pxl}

\begin{frame}[label=p2]{Simulation: count the rotations}
  \setbeamercovered{transparent=32}
  \vspace{0.5em}
  \begin{columns}[T,onlytextwidth]
  \begin{column}{0.60\textwidth}
  \begin{center}
  \begin{tikzpicture}[x=1cm,y=1cm,scale=0.84,every node/.style={transform shape}]
    \setpinscale{1}
    % column = sorted position (5,10,20,30,40,60), row = depth
    \foreach \cc in {0,...,5}{\foreach \dd in {0,1,2,3}{
      \coordinate (c\cc-\dd) at ({(\cc-2.5)*1.25},{-\dd*1.20});}}

    % Fixed extent: the tree must not resize between beats. The top is set by
    % the turn arrow and its label, the bottom by the ghost of 5.
    \path (-3.70,-3.98) rectangle (3.70,1.50);

    % ---------------- the tree as it stands once 5 has arrived -------------
    \only<1-4>{
      \draw[avledge] (c4-0) -- (c2-1) -- (c1-2) -- (c0-3);
      \draw[avledge] (c4-0) -- (c5-1);
      \draw[avledge] (c2-1) -- (c3-2);
    }
    % the walk itself, drawn on the links it climbs
    \only<2-4>{\draw[draw=avlptwo,line width=2\pxl] (c0-3) -- (c1-2) -- (c2-1);}
    \only<3-4>{\draw[draw=avlptwo,line width=2\pxl] (c2-1) -- (c4-0);}

    \only<1>{
      \node[avlkey]     at (c4-0) {40}; \node[avlkey] at (c2-1) {20};
      \node[avlkey]     at (c5-1) {60}; \node[avlkey] at (c1-2) {10};
      \node[avlkey]     at (c3-2) {30}; \node[avlkeygold] at (c0-3) {5};
      \node[avlannotfib,anchor=west] at (-2.62,-3.60) {new};
    }
    % Walk beats. A node that has been checked and passed keeps its green, so
    % the audience can see how much of the walk found nothing.
    \only<2-4>{
      \node[avlkeyok] at (c2-1) {20}; \node[avlkey] at (c5-1) {60};
      \node[avlkeyok] at (c1-2) {10}; \node[avlkey] at (c3-2) {30};
      \node[avlkeygold] at (c0-3) {5};
      \node[avlannotok,anchor=east] at (-2.38,-2.40) {$\mathrm{bf} = +1$};
      \node[avlannotok,anchor=east] at (-1.13,-1.20) {$\mathrm{bf} = +1$};
    }
    \only<2>{\node[avlkey] at (c4-0) {40};}
    \only<3-4>{\node[avlkeyalert] at (c4-0) {40};}
    \only<3>{
      \draw[avlalert,line width=2\pxl] (c4-0) circle (0.42);
      \node[avlannotalert,anchor=west] at (2.45,0.18) {$\mathrm{bf} = +2$};}
    \only<4>{
      \pivotpin{c4-0}\rotright{c4-0}{1.30}
      \node[color=avlptwo,font=\avlfAnnot\wmed,anchor=east] at (0.40,1.28)
        {LL: turn right at 40};}

    % ---------------- the turn has landed, afterimages still up ------------
    \only<5>{
      \ghostnode{1.875}{0}{40}       \ghostpath{1.875,-0.35}{1.875,-0.85}
      \ghostnode{-0.625}{-1.20}{20}  \ghostpath{-0.625,-0.85}{-0.625,-0.35}
      \ghostnode{-3.125}{-3.60}{5}   \ghostpath{-3.125,-3.25}{-3.125,-2.75}
      \draw[avlgoldedge,line width=1.1pt] (c1-1) -- (c2-0) -- (c4-1);
      \draw[avlgoldedge,line width=1.1pt] (c0-2) -- (c1-1);
      \draw[avlgoldedge,line width=1.1pt] (c3-2) -- (c4-1) -- (c5-2);
      \node[avlkeygold] at (c2-0) {20}; \node[avlkeygold] at (c1-1) {10};
      \node[avlkeygold] at (c4-1) {40}; \node[avlkeygold] at (c0-2) {5};
      \node[avlkeygold] at (c3-2) {30}; \node[avlkeygold] at (c5-2) {60};
    }
    % ---------------- verdict: same tree, afterimages cleared, green -------
    % Nothing moves between beat 5 and beat 6. The only change is that the
    % tree stops being a thing in motion and becomes a legal result.
    \only<6>{
      \draw[draw=avlbalance!70,line width=1.6\pxl] (c1-1) -- (c2-0) -- (c4-1);
      \draw[draw=avlbalance!70,line width=1.6\pxl] (c0-2) -- (c1-1);
      \draw[draw=avlbalance!70,line width=1.6\pxl] (c3-2) -- (c4-1) -- (c5-2);
      \node[avlkeyok] at (c2-0) {20}; \node[avlkeyok] at (c1-1) {10};
      \node[avlkeyok] at (c4-1) {40}; \node[avlkeyok] at (c0-2) {5};
      \node[avlkeyok] at (c3-2) {30}; \node[avlkeyok] at (c5-2) {60};
      \node[font=\avlfLead\wbold,text=avlbalance,anchor=north] at (0,-3.05)
        {Legal again. One rotation.};
    }
  \end{tikzpicture}
  \end{center}
  \end{column}

  % ---------------- the algorithm, tracked beside the tree ----------------
  \begin{column}{0.36\textwidth}
  \vspace{1.1em}
  \insstep{1}{Insert like a normal BST}
  \insstep{2-3}{Walk back up to the root}
  \insstep{4-5}{Rotate when we find a bad node}
  \insstep[6]{6}{Done}
  \end{column}
  \end{columns}

  \note{A tree they have not seen. Click: 5 arrives, ordinary BST insert, no
  balance work at all yet. Click: now walk back up. 10 is fine, 20 is fine.
  Say out loud that this is the common case, the walk usually finds nothing.
  Click: the root is the one that broke, plus two. Click: left child, left
  side, so it is LL. Pin 40, turn clockwise. Click: the turn has landed, and
  the faded circles say where 40, 20 and 5 came from. Click: legal. Count it
  with them, one rotation. Then ask whether that was luck.}
\end{frame}

% ------------------------------------------------------------------ the catch
%  ONE box on this slide, and it is the answer. The questions are set as plain
%  type with real air around them, so the only bordered thing on the ground is
%  the thing the audience is waiting for. The hierarchy is carried by size --
%  17 pt questions against a 25 pt answer -- and by colour: ink for the
%  question anyone would ask, part hue for the one that presses, green for the
%  verdict. A second card around the questions would have flattened all of it
%  back into two grey rectangles.
\begin{frame}[label=p2]{The catch}
  \vspace{0.6em}
  \begin{center}
    % The baselineskip in force here is \avlfBody's 14.4 pt, not \avlfStmt's:
    % the size command is scoped inside the braces, so the \\ that follows it
    % still breaks at body leading. That is why these gaps look outsized.
    {\avlfStmt\wmed\color{avlink}One rotation and we are done?}\\[72\pxl]
    \uncover<2->{{\avlfStmt\wmed\color{avlptwo}Always the case?}}\\[78\pxl]
    % Fixed slot: the answer must not push the questions up the slide when it
    % arrives. Sized to the panel plus its trailing line.
    \begin{minipage}[t][174\pxl][t]{\textwidth}
    \centering
    \only<3->{%
      \begin{avlpanelh}[0.46\textwidth]{avlbalance}
        {\avlfBig\wbold\color{avlbalance}Yes, it is.}
      \end{avlpanelh}\\[18\pxl]
      {\avlfSec\color{avlmuted}At most one rotation per insert.}}
    \end{minipage}
  \end{center}

  \note{Slow down here. One rotation. Was that this tree, or is it a
  guarantee? Let them sit with it for a beat. Click. Ask it again, out loud:
  always? Click. Yes, always, at most one rotation per insert, and that is not
  an observation, it is provable. The next slide is the proof in one picture.}
\end{frame}

% ------------------------------------------------------------------ why so
%  The height argument, drawn rather than asserted: the insert lifts one
%  subtree by exactly one, the rotation puts that one back, and an ancestor
%  that sees the height it saw before has nothing left to repair.
\begin{frame}[label=p2]{Why so?}
  \vspace{-0.1em}
  \begin{center}
  \begin{tikzpicture}[x=1cm,y=1cm,scale=0.75,every node/.style={transform shape},
      par/.style={circle,draw=avlink!70,fill=white!96!avlground,
                  line width=1.5\pxl,inner sep=0pt,minimum size=30\pxl},
      whycap/.style={font=\avlfAnnot,text=avlmuted,anchor=north,align=center}]
    \path (-6.30,-1.80) rectangle (6.30,1.45);

    % ---- before the insert : height h ----
    \begin{scope}[shift={(-4.60,0)}]
      \node[par] at (0,1.20) {};
      \draw[avledge] (0,1.20) -- (0,0.85);
      \draw[draw=avlfib!55,fill=avlfib!8!avlground,line width=1.6\pxl]
        (0,0.85) -- (-1.00,-0.85) -- (1.00,-0.85) -- cycle;
      \heightbrace{-1.32}{0.85}{-0.85}{$h$}{avlfaint}
      \node[whycap] at (0,-1.05) {before the insert};
    \end{scope}

    % ---- after the insert : one taller ----
    \begin{scope}
      \node[par] at (0,1.20) {};
      \draw[avledge] (0,1.20) -- (0,0.85);
      \draw[draw=avlalert!55,fill=avlalert!8!avlground,line width=1.6\pxl]
        (0,0.85) -- (-1.15,-1.25) -- (1.15,-1.25) -- cycle;
      \heightbrace{-1.47}{0.85}{-1.25}{$h+1$}{avlalert}
      \node[whycap] at (0,-1.45) {after the insert};
    \end{scope}

    % ---- after one rotation : height h again ----
    % Pushed a little further right than the mirror position: this brace sits
    % on the same y as the "rotate" arrow, and its rotated label needs clear
    % air between itself and the arrowhead.
    \begin{scope}[shift={(4.80,0)}]
      % Whole state hidden until the rotation beat -- the parent stub on its
      % own reads as a stray lollipop with no subtree under it.
      \only<2->{
        \node[par] at (0,1.20) {};
        \draw[avledge] (0,1.20) -- (0,0.85);
        \draw[draw=avlbalance!60,fill=avlbalance!8!avlground,line width=1.6\pxl]
          (0,0.85) -- (-1.00,-0.85) -- (1.00,-0.85) -- cycle;
        \heightbrace{-1.32}{0.85}{-0.85}{$h$}{avlbalance}
        \node[whycap] at (0,-1.05) {after one rotation};}
    \end{scope}

    \draw[avlgoldarrow] (-3.35,-0.10) -- (-1.95,-0.10)
      node[midway,above=2\pxl,font=\avlfAnnot,text=avlfib] {insert};
    \only<2->{
      \draw[avlarc] (1.70,-0.10) -- (2.90,-0.10)
        node[midway,above=2\pxl,font=\avlfAnnot,text=avlptwo] {rotate};}
  \end{tikzpicture}
  \end{center}

  \vspace{0.2em}
  \begin{center}
  % HARD BUDGET. Below the figure there are 2.45 cm, and a fixed-height
  % minipage hides an overrun from the compiler: a wrapped line here spills
  % silently through the footer rule instead of raising an overfull warning.
  % Both sentences and the panel line must each stay on ONE line, which is
  % roughly 72 characters at \avlfSec and roughly 46 at \avlfLead.
  \begin{minipage}[t][196\pxl][t]{0.94\textwidth}
  \centering
    {\avlfSec\color{avlmuted}An insert adds
     \viol{at most one level} to a subtree.}\\[0.5em]
    \uncover<2->{{\avlfSec\color{avlmuted}The rotation takes that level back,
      so \ok{nothing above the node changes}.}}\\[0.6em]
    \only<3->{%
      \begin{avlpanel}[0.94\linewidth]
        {\avlfLead\color{avlfib}\wbold Insertion is $O(\log n)$: a walk and
         one turn.}
      \end{avlpanel}}
  \end{minipage}
  \end{center}

  \note{Here is the whole argument in one picture. An insert adds one node, so
  it can lift a subtree by at most one level. Click. The rotation gives that
  level straight back: the subtree comes out at exactly the height it had
  before the insert, so the parent, and everything above the parent, sees no
  change at all and has nothing to fix. That is why the repair stops at the
  node we rotated. Click. And the cost: log n down, log n back up, one
  constant-time rotation. Insertion is logarithmic, guaranteed.}
\end{frame}

% ===========================================================================
%  PART 3 --- Deletion, the worst case simulation, and the Fibonacci payoff
%  Flow: deletion rule -> gentle example -> cascade -> Fibonacci tie-back
%        -> rotation bound -> applications -> closing
%
%  Every number is simulated ground truth:
%    worst-case deletion rotations = floor(h/2) exactly, h = 2..14
%    worst-case insertion rotations = 1 at every height tested
%    h=7 tree: deleting the leaf of in-order rank 33 (pre-order f34) fires
%      3 rotations, at ranks 32, 29, 21 (pre f31, f23, f2), depths 3 -> 2 -> 1
%    per-layer subtree sizes: 54 33 20 12 7 4 2 1, and the rotation sites are
%      the T_2, T_4, T_6 layers -- every second one
%  The cascade geometry is regenerated by verify/cascade_gen.py, which asserts
%  all of the above before it writes tikz/cascade.tex. Re-run it after any
%  change to the tree, and never hand-edit the generated file.
% ===========================================================================

% Generated geometry for the deletion cascade below. Loaded HERE rather than
% from avl.tex on purpose, so this part is self-contained: avl.tex needs NO
% change for this to work. These are plain \def's, so loading them inside the
% document body is fine.
%% GENERATED by verify/cascade_gen.py -- do not edit by hand.
%% Worst-case deletion cascade on the minimal height-7 AVL tree.
%% Delete leaf of in-order rank 33 (pre-order f34). Three rotations fire,
% at ranks 32, 29, 21 (pre f31, f23, f2) at depths 3, 2, 1.
%% x of every node is its ORIGINAL in-order rank, held fixed for the whole
%% cascade, so nodes only ever move VERTICALLY -- the same rule Part 2 states.
%% Rank 33 is absent from every state after the first: the node is really gone.

%% --- initial: 54 nodes, height 7
\def\cascNZero{1/7,2/6,3/5,4/6,5/4,6/6,7/5,8/3,9/6,10/5,11/4,12/5,13/2,14/6,15/5,16/4,17/5,18/3,19/5,20/4,21/1,22/6,23/5,24/4,25/5,26/3,27/5,28/4,29/2,30/5,31/4,32/3,33/4,34/0,35/6,36/5,37/4,38/5,39/3,40/5,41/4,42/2,43/5,44/4,45/3,46/4,47/1,48/5,49/4,50/3,51/4,52/2,53/4,54/3}
\def\cascEZero{34/21,21/13,13/8,8/5,5/3,3/2,2/1,3/4,5/7,7/6,8/11,11/10,10/9,11/12,13/18,18/16,16/15,15/14,16/17,18/20,20/19,21/29,29/26,26/24,24/23,23/22,24/25,26/28,28/27,29/32,32/31,31/30,32/33,34/47,47/42,42/39,39/37,37/36,36/35,37/38,39/41,41/40,42/45,45/44,44/43,45/46,47/52,52/50,50/49,49/48,50/51,52/54,54/53}
\def\cascHZero{7}

%% --- before rot 1: 53 nodes, height 7, |bf|=2 at rank 32 (+2)
\def\cascNOne{1/7,2/6,3/5,4/6,5/4,6/6,7/5,8/3,9/6,10/5,11/4,12/5,13/2,14/6,15/5,16/4,17/5,18/3,19/5,20/4,21/1,22/6,23/5,24/4,25/5,26/3,27/5,28/4,29/2,30/5,31/4,32/3,34/0,35/6,36/5,37/4,38/5,39/3,40/5,41/4,42/2,43/5,44/4,45/3,46/4,47/1,48/5,49/4,50/3,51/4,52/2,53/4,54/3}
\def\cascEOne{34/21,21/13,13/8,8/5,5/3,3/2,2/1,3/4,5/7,7/6,8/11,11/10,10/9,11/12,13/18,18/16,16/15,15/14,16/17,18/20,20/19,21/29,29/26,26/24,24/23,23/22,24/25,26/28,28/27,29/32,32/31,31/30,34/47,47/42,42/39,39/37,37/36,36/35,37/38,39/41,41/40,42/45,45/44,44/43,45/46,47/52,52/50,50/49,49/48,50/51,52/54,54/53}
\def\cascHOne{7}
\def\cascHotOne{32}
\def\cascBfOne{+2}

%% --- before rot 2: 53 nodes, height 7, |bf|=2 at rank 29 (+2)
\def\cascNTwo{1/7,2/6,3/5,4/6,5/4,6/6,7/5,8/3,9/6,10/5,11/4,12/5,13/2,14/6,15/5,16/4,17/5,18/3,19/5,20/4,21/1,22/6,23/5,24/4,25/5,26/3,27/5,28/4,29/2,30/4,31/3,32/4,34/0,35/6,36/5,37/4,38/5,39/3,40/5,41/4,42/2,43/5,44/4,45/3,46/4,47/1,48/5,49/4,50/3,51/4,52/2,53/4,54/3}
\def\cascETwo{34/21,21/13,13/8,8/5,5/3,3/2,2/1,3/4,5/7,7/6,8/11,11/10,10/9,11/12,13/18,18/16,16/15,15/14,16/17,18/20,20/19,21/29,29/26,26/24,24/23,23/22,24/25,26/28,28/27,29/31,31/30,31/32,34/47,47/42,42/39,39/37,37/36,36/35,37/38,39/41,41/40,42/45,45/44,44/43,45/46,47/52,52/50,50/49,49/48,50/51,52/54,54/53}
\def\cascHTwo{7}
\def\cascHotTwo{29}
\def\cascBfTwo{+2}

%% --- before rot 3: 53 nodes, height 7, |bf|=2 at rank 21 (+2)
\def\cascNThree{1/7,2/6,3/5,4/6,5/4,6/6,7/5,8/3,9/6,10/5,11/4,12/5,13/2,14/6,15/5,16/4,17/5,18/3,19/5,20/4,21/1,22/5,23/4,24/3,25/4,26/2,27/5,28/4,29/3,30/5,31/4,32/5,34/0,35/6,36/5,37/4,38/5,39/3,40/5,41/4,42/2,43/5,44/4,45/3,46/4,47/1,48/5,49/4,50/3,51/4,52/2,53/4,54/3}
\def\cascEThree{34/21,21/13,13/8,8/5,5/3,3/2,2/1,3/4,5/7,7/6,8/11,11/10,10/9,11/12,13/18,18/16,16/15,15/14,16/17,18/20,20/19,21/26,26/24,24/23,23/22,24/25,26/29,29/28,28/27,29/31,31/30,31/32,34/47,47/42,42/39,39/37,37/36,36/35,37/38,39/41,41/40,42/45,45/44,44/43,45/46,47/52,52/50,50/49,49/48,50/51,52/54,54/53}
\def\cascHThree{7}
\def\cascHotThree{21}
\def\cascBfThree{+2}

%% --- final: 53 nodes, height 6
\def\cascNFour{1/6,2/5,3/4,4/5,5/3,6/5,7/4,8/2,9/5,10/4,11/3,12/4,13/1,14/6,15/5,16/4,17/5,18/3,19/5,20/4,21/2,22/6,23/5,24/4,25/5,26/3,27/6,28/5,29/4,30/6,31/5,32/6,34/0,35/6,36/5,37/4,38/5,39/3,40/5,41/4,42/2,43/5,44/4,45/3,46/4,47/1,48/5,49/4,50/3,51/4,52/2,53/4,54/3}
\def\cascEFour{34/13,13/8,8/5,5/3,3/2,2/1,3/4,5/7,7/6,8/11,11/10,10/9,11/12,13/21,21/18,18/16,16/15,15/14,16/17,18/20,20/19,21/26,26/24,24/23,23/22,24/25,26/29,29/28,28/27,29/31,31/30,31/32,34/47,47/42,42/39,39/37,37/36,36/35,37/38,39/41,41/40,42/45,45/44,44/43,45/46,47/52,52/50,50/49,49/48,50/51,52/54,54/53}
\def\cascHFour{6}

\def\casctarget{33}

\avldivider{3}{Deletion}{The expensive one.}{avlpthree}{avlpthreeA}{avlpthreeB}{avlpthreeC}

\setavlpart{3}{Part 3 --- Deletion}{avlpthree}

\begin{frame}[label=p3]{Deletion: the rule}
  \vspace{0.5em}
  \begin{center}
  \begin{tikzpicture}[x=1cm,y=1cm,
      stepbox/.style={draw=avlpthree!45!avlground,fill=avlpthree!8!avlground,
                   line width=1.5\pxl,rounded corners=6pt,
                   minimum width=3.6cm,minimum height=1.15cm,
                   font=\avlfSec,text=avlink},
      bad/.style ={draw=avlalert!45!avlground,fill=avlalert!8!avlground,
                   line width=1.5\pxl,rounded corners=6pt,
                   minimum width=3.6cm,minimum height=1.15cm,
                   font=\avlfSec,text=avlink}]
    \node[stepbox] (s1) at (-4.5,0)
      {\begin{minipage}{3.2cm}\centering Remove like a normal BST\end{minipage}};
    \node[stepbox] (s2) at (0,0)
      {\begin{minipage}{3.2cm}\centering Walk back up to the root\end{minipage}};
    \node[bad] (s3) at (4.5,0)
      {\begin{minipage}{3.2cm}\centering Rotate at \strong{every} bad node\end{minipage}};
    \draw[avlarrow] (s1) -- (s2);
    \draw[avlarrow] (s2) -- (s3);
  \end{tikzpicture}
  \end{center}

  \vspace{1.0em}
  \begin{center}
    {\avlfLead\color{avlbody}The first two steps are the same as insertion.}\\[0.7em]
    \begin{avlpanelh}[0.62\textwidth]{avlalert}
      {\avlfLead\viol{\wbold The third one is not.}}
    \end{avlpanelh}\\[0.8em]
    {\avlfSec\color{avlmuted}Insertion stops after one single or double repair.
     Deletion may need repairs higher up.}
  \end{center}

  \note{The first two steps look identical to insertion. The difference is the
  last word. Insertion fixes one node and stops. Deletion may have to fix every
  node on the way up.}
\end{frame}

\begin{frame}[label=p3]{Deletion: same rule, new twist}
  \vspace{-0.4em}
  \begin{center}
  \begin{tikzpicture}[x=1cm,y=1cm]
    \path (-5.6,-2.0) rectangle (5.6,1.05);
    \only<1-2>{
      \node[avlannot] at (-3.2,0.65) {delete 10};
      \draw[avledge] (-4.3,-0.8) -- (-3.2,0) -- (-2.1,-0.8);
      \node[avlkey] at (-3.2,0) {20};
      \node[avlkeyalert] at (-4.3,-0.8) {10};
      \node[avlkey] at (-2.1,-0.8) {30};
    }
    \only<2>{
      \draw[avlarrow] (-0.6,-0.45) -- (0.6,-0.45);
      \ghostnode{1.8}{-0.8}{10}
      \draw[avlghostline] (2.9,0) -- (1.8,-0.8);
      \draw[avledge] (2.9,0) -- (4.0,-0.8);
      \node[avlkeyok] at (2.9,0) {20};
      \node[avlkeyok] at (4.0,-0.8) {30};
      \node[avlannotok] at (2.9,0.65) {$\mathrm{bf}(20)=-1$};
      \node[avlannotok] at (2.9,-1.65) {still legal; no rotation};
    }
    \only<3-6>{
      \node[avlannot,anchor=west,align=left,text width=1.8cm] at (-5.3,0.65) {a taller tree\\delete 30};
      \only<3>{
        \draw[avledge] (-2.1,0) -- (-1.0,-0.8);
        \node[avlkeyalert] at (-1.0,-0.8) {30};
      }
      \only<4-6>{
        \ghostnode{-1.0}{-0.8}{30}
        \draw[avlghostline] (-2.1,0) -- (-1.0,-0.8);
      }
      \only<5-6>{\rotright{-2.1,0}{0.9}}
      \draw[avledge] (-2.1,0) -- (-3.2,-0.8) -- (-4.3,-1.6);
      \only<3>{\node[avlkey] at (-2.1,0) {20};}
      \only<4-6>{\node[avlkeyalert] at (-2.1,0) {20};}
      \node[avlkey] at (-3.2,-0.8) {10};
      \node[avlkeygold] at (-4.3,-1.6) {5};
      \only<5-6>{\pivotpin[avlalert]{-2.1,0}}
      \only<4-6>{
        \node[avlannotalert,anchor=east] at (-2.6,-0.1) {$\mathrm{bf}=+2$};
      }
    }
    \only<6>{
      \draw[avlarrow] (-0.4,-0.45) -- (0.8,-0.45);
      \ghostnode{4.0}{0}{20}
      \ghostpath{4.0,-0.3}{4.0,-0.5}
      \ghostnode{2.9}{-0.8}{10}
      \ghostpath{2.9,-0.5}{2.9,-0.3}
      \ghostnode{1.8}{-1.6}{5}
      \ghostpath{1.8,-1.3}{1.8,-1.1}
      \draw[draw=avlbalance!70,line width=1.6\pxl]
        (1.8,-0.8) -- (2.9,0) -- (4.0,-0.8);
      \node[avlkeyok] at (2.9,0) {10};
      \node[avlkeyok] at (1.8,-0.8) {5};
      \node[avlkeyok] at (4.0,-0.8) {20};
      \node[avlannotok] at (2.9,0.65) {one right rotation};
    }
  \end{tikzpicture}

  \begin{minipage}[t][74\pxl][t]{\textwidth}
    \centering\avlfSec\color{avlmuted}
    \only<1>{Start with 20 above 10 and 30.
      Every node is balanced.\par}
    \only<2>{Delete 10. 20 just loses a child --- $\mathrm{bf}(20)=-1$,
      still legal. Nothing to fix.\par}
    \only<3>{A fresh, taller tree. $\mathrm{bf}(20)=+1$, still legal.\\
      This time, delete 30.\par}
    \only<4>{30 is gone. Now $\mathrm{bf}(20)=+2$.\\
      The same alarm from Part 2 goes off.\par}
    \only<5>{Pin 20 and turn clockwise.\\
      One right rotation, just as in Part 2.\par}
    \only<6>{10 becomes the root. The tree is balanced again.\\
      One rotation --- so far, no different from insertion.\par}
  \end{minipage}
  \begin{avlpanelh}[0.94\textwidth]{avlfib}
    \avlfSec\color{avlbody}
    The rule reads the same as insertion.\\
    The difference shows up somewhere else.\par
  \end{avlpanelh}
  \end{center}
  \note{Six clicks, with the before and after separated.
  First, show the balanced three-node tree. Second, delete 10 and reveal the
  result on the right. The balance factor of 20 becomes minus one; no rotation.
  Third, show the fresh taller tree before deleting anything: 20 has children
  10 and 30, and 10 has left child 5. Fourth, delete 30; the balance factor of
  20 becomes plus two. Fifth, introduce the pin and clockwise arc at 20.
  Sixth, reveal the repaired tree on the right: 10 above 5 and 20, with ghosts
  at the old positions. The taller tree is a new example, not a continuation
  of the first deletion.}
\end{frame}

\begin{frame}[label=p3]{From one rotation to a cascade}
  \begin{center}
    {\avlfStmt\wbold\color{avlink}When does a rotation\\
     propagate toward the root?\par}

    \vspace{1.4em}
    {\avlfBody\color{avlmuted}We fixed one node. Is the parent still balanced?\par}
  \end{center}
  \note{Pause here and let the question land. The small example needed only
  one rotation. What could make a repair continue upward instead of stopping?
  Invite the audience to think about what changes for the parent, then show
  the worst-case tree and let the cascade reveal the answer.}
\end{frame}

% ------------------------------------------------------------ WORST CASE SIM
%  The cascade, run for real. Every overlay below is a DIFFERENT tree: the node
%  is actually removed, and each rotation actually restructures the picture.
%  All five states come from verify/cascade_gen.py driving a real AVL
%  implementation; the geometry is read in from tikz/cascade.tex.
%
%  Layout rule (the same one Part 2 states): a node's x is its ORIGINAL
%  in-order rank and never changes. Only depth changes. So a rotation moves
%  nodes straight up or down, the deleted node leaves a permanent gap at its
%  own x, and the height drop at the end is a visible change in the picture
%  rather than a claim in the caption.
\newcommand{\cascdx}{0.196}
\newcommand{\cascdy}{0.50}
\newcommand{\cascrad}{1.5pt}

% \casctree{<nodes>}{<edges>} -- edges behind, discs on top
\newcommand{\casctree}[2]{%
  \foreach \r/\d in #1 {\coordinate (c-\r) at (\r*\cascdx,-\d*\cascdy);}%
  \foreach \a/\b in #2 {\draw[avlgoldedge] (c-\a) -- (c-\b);}%
  \foreach \r/\d in #1 {\fill[avlfib] (c-\r) circle (\cascrad);}}

% The node that is gone. It used to be marked with a dashed ring held at the
% x-slot it would never occupy again; that ring read as a stray artefact once
% the tree had restructured around it, so the slot is now simply left empty.
\newcommand{\cascgap}{}

% \cascbreak{<rank>}{<bf>} -- the node that is now illegal
\newcommand{\cascbreak}[2]{%
  \draw[draw=avlalert, line width=1.2pt] (c-#1) circle (4.2pt);
  % The label goes ABOVE THE ROOT ROW, the one band of this figure that is
  % guaranteed empty at every beat, and a hairline leader ties it back to the
  % ring. Anywhere inside the canopy its filled box lands on some edge and
  % chops it -- which is what it was doing next to the node.
  \draw[avlalert!50, line width=0.5pt]
    ($(c-#1)+(0,4.6pt)$) -- ({#1*\cascdx},0.10);
  \node[text=avlalert, font=\scriptsize\bfseries, anchor=south,
        fill=avlground, inner sep=1.8pt, rounded corners=1pt]
    at ({#1*\cascdx},0.10) {$\mathrm{bf}=#2$};}

% height ruler down the left edge; #1 = height, #2 = colour
\newcommand{\cascruler}[2]{%
  \begin{scope}[xshift=-0.55cm]
    % mirror: the brace must open toward the tree it is measuring
    \draw[decorate,decoration={brace,amplitude=4pt,mirror},#2,line width=0.8pt]
      (0,0.10) -- (0,{-#1*\cascdy-0.10});
    \node[color=#2,font=\scriptsize,anchor=east,rotate=90]
      at (-0.26,{-#1*\cascdy/2}) {height #1};
  \end{scope}}

\begin{frame}[label=p3]{Worst case: one leaf out, and it does not stop}
  \vspace{-0.5em}
  \begin{center}
  \begin{tikzpicture}[x=1cm,y=1cm,scale=0.90,every node/.append style={transform shape}]
    % fixed extent: the picture must not resize between clicks
    \path (-1.5,-3.70) rectangle (11.2,0.60);

    \only<1>{\casctree{\cascNZero}{\cascEZero}%
             \draw[draw=avlalert,line width=1.2pt]
               (\casctarget*\cascdx,-4*\cascdy) circle (3.6pt);
             \cascruler{\cascHZero}{avlfib}}
    \only<2-3>{\casctree{\cascNOne}{\cascEOne}\cascgap
             \cascbreak{\cascHotOne}{\cascBfOne}\cascruler{\cascHOne}{avlfib}}
    \only<4>{\casctree{\cascNTwo}{\cascETwo}\cascgap
             \cascbreak{\cascHotTwo}{\cascBfTwo}\cascruler{\cascHTwo}{avlfib}}
    \only<5>{\casctree{\cascNThree}{\cascEThree}\cascgap
             \cascbreak{\cascHotThree}{\cascBfThree}\cascruler{\cascHThree}{avlfib}}
    \only<6->{\casctree{\cascNFour}{\cascEFour}\cascgap
             \cascruler{\cascHFour}{avlpthree}}

    \only<3>{
      \begin{scope}[overlay]
        \setpinscale{0.22}
        \rotright{c-32}{0.42}
        \pivotpin[avlalert]{c-32}
        \node[avlannotalert,anchor=north west,align=left,text width=3.5cm]
          (shorterNote) at (7.5,-2.65)
          {Unlike insertion: one level shorter. The parent can break too.};
        \draw[avlalert!55,line width=1\pxl]
          ($(c-32)+(0.38,-0.05)$) -- (7.35,-2.48) -- (shorterNote.north west);
      \end{scope}
    }
  \end{tikzpicture}
  \end{center}

  % Caption block reserves its own height: six beats, one of them two lines.
  \vspace{-0.7em}
  \begin{center}
  \begin{minipage}[t][86\pxl][t]{\textwidth}\centering\avlfSec
    \only<1>{{\color{avlmuted}54 nodes, height 7, no spare room anywhere.
              Delete the \viol{marked leaf}.}}
    \only<2-3>{{\color{avlmuted}Gone. Its grandparent is now
              \viol{$\mathrm{bf}=+2$} --- illegal.
              {\color{avlpthree}\wmed Rotate.}}}
    \only<4>{{\color{avlmuted}Fixed --- but that subtree got \emph{shorter},
              so its own parent is now \viol{$+2$}.
              {\color{avlpthree}\wmed Rotate again.}}}
    \only<5>{{\color{avlmuted}And again. This one is a child of the root.
              {\color{avlpthree}\wmed Rotate a third time.}}}
    \only<6->{{\avlfLead\viol{\wbold 3 rotations}\quad\quietc{$\cdot$}\quad
               {\color{avlbody}54 $\rightarrow$ 53 nodes}
               \quad\quietc{$\cdot$}\quad
               {\color{avlbody}height \textbf{7} $\rightarrow$}
               \color{avlpthree}\textbf{6}}\\[0.35em]
              {\color{avlmuted}The height only fell when the cascade reached
               the top --- and one delete paid for all of it.}}
  \end{minipage}
  \end{center}

  \vspace{-0.1em}
  {\avlfAnnot\color{avlmuted}\centering
   Note: this cascade is left-heavy ($\mathrm{bf}=+2$); right-heavy is its mirror.\\
   Height loss can propagate after any Part 2 repair: LL, RR, LR, or RL.\par}

  \note{Keep all five tree states. The extra beat pauses at the first rotation
  trigger so you can point to its pin and arc. This rotation shortens the
  subtree by one --- that is the difference from insertion. A shorter subtree
  can make the parent illegal too. Then continue the original cascade:
  fix the first site, fix its parent, and fix the child of the root.
  The whole tree drops from height seven to six after three rotations.
  This example propagates height loss; not every deletion rotation does.}
\end{frame}

\newcommand{\avldeletionlayers}{%
  \begin{tikzpicture}[x=1cm,y=1cm,
      avldot/.style={circle,draw=avlfib,fill=avlfib,inner sep=0pt,
                     minimum size=5.8\pxl},
      avledge/.style={draw=avlfib!75,line width=1.3\pxl}]
    \setavlscale{0.42}{0.33}
    \resetavl
    \drawfib{7}{0}{0}

    % The root's tag would sit under the title rule if it went left like the
    % rest, so it is the one that leans right instead.
    \node[avlannotfib,anchor=west]
      at ($(f1)+(0.22,0.10)$) {$T_{7}$\,:\,\textbf{54}};
    \foreach \i/\k/\N in {2/6/33, 3/5/20, 4/4/12,
                          5/3/7, 6/2/4, 7/1/2, 8/0/1} {
      \draw[avlfib!45,line width=1\pxl]
        ($(f\i)+(-0.14,0)$) -- ($(f\i)+(-0.30,0)$);
      \node[avlannotfib,anchor=east]
        at ($(f\i)+(-0.32,0.01)$) {$T_{\k}$\,:\,\textbf{\N}};
    }
  \end{tikzpicture}
}

\begin{frame}[t,label=p3]{A Fibonacci tree at every level}
  \vspace{-0.2em}
  \begin{center}
    \avldeletionlayers

    \vspace{0.4em}
    {\avlfSec\color{avlmuted}Each subtree is a smaller \strong{Fibonacci tree}.\\
     It uses the fewest nodes possible for its height.\par}
    \vspace{0.5em}
    {\avlfSec\color{avlfib}$54,\ 33,\ 20,\ 12,\ 7,\ 4,\ 2,\ 1$
     \quad\quietc{---}\quad $N(7)$ down to $N(0)$.\par}
  \end{center}
  \note{The cascade tree is a Fibonacci tree: a minimal AVL tree at every
  level. Each subtree combines the smallest trees of the two preceding
  heights. Its node count follows N(h) = 1 + N(h-1) + N(h-2), so the counts
  are Fibonacci numbers minus one. Read down the left edge: 54, 33, 20, 12,
  7, 4, 2, 1. That is N(7) down to N(0). Next, connect that minimum node
  count to the height loss that drives the cascade.}
\end{frame}

\begin{frame}[t,label=p3]{Why the Fibonacci tree cascades}
  \vspace{-0.2em}
  \begin{center}
    \avldeletionlayers

    \vspace{0.3em}
    {\avlfSec\color{avlmuted}Fibonacci numbers in disguise:
     \fib{$54 = F(10)-1$}.\par}
    \vspace{0.3em}
    {\avlfSec\color{avlmuted}Height 7 needs at least 54 nodes.\\
     Delete one leaf: 53 nodes cannot keep that height.\par}
    \vspace{0.5em}
    \begin{avlpanelh}[0.94\textwidth]{avlpthree}
      \avlfSec\color{avlpthree}No slack anywhere. That's why the fix cascades.\par
    \end{avlpanelh}
  \end{center}
  \note{Same diagram, same thought. The subtree counts are Fibonacci numbers
  minus one: N(h) equals F(h+3) minus one. In particular, 54 is F(10) minus one.
  This tree has no slack. Removing a leaf must reduce its height after repair:
  a height-seven AVL tree needs at least 54 nodes, and only 53 remain.
  Our chosen leaf forces that repair to cascade. Not every leaf must cause
  three rotations; this is the deliberately chosen worst-case deletion.}
\end{frame}

\begin{frame}[label=p3]{Insertion costs 1. Deletion costs $\lfloor h/2 \rfloor$.}
  % Fixed two-column layout: the plot is drawn at its true size (640 px wide)
  % and is never scaled to fit, so its tick and axis labels stay legible.
  % No grid -- the ticks carry the scale, and a guide line under the series
  % labels is the one thing that kept colliding with them.
  \vspace{0.2em}
  \begin{columns}[c,onlytextwidth]
    \begin{column}{7.9cm}
      \begin{tikzpicture}
        \begin{axis}[
          width=7.3cm, height=4.45cm, axis lines=left, scale only axis,
          xlabel={height of the tree}, ylabel={rotations, worst case},
          xlabel style={font=\avlfAnnot,color=avlfaint},
          ylabel style={font=\avlfAnnot,color=avlfaint},
          tick label style={font=\avlfAnnot,color=avlfaint},
          axis line style={color=avlfaint,line width=1\pxl},
          tick style={color=avlfaint,line width=1\pxl},
          xtick={2,4,6,8,10,12,14}, ytick={0,2,4,6,8},
          ymin=0, ymax=8, xmin=2, xmax=14,
        ]
          % the 13 measured points, verify/ ground truth
          \addplot[color=avlalert,line width=2.4\pxl,mark=*,mark size=2\pxl,
                   mark options={fill=avlalert,draw=avlalert}]
            coordinates {(2,1)(3,1)(4,2)(5,2)(6,3)(7,3)(8,4)(9,4)(10,5)(11,5)
                         (12,6)(13,6)(14,7)};
          \addplot[color=avlfib,line width=2.4\pxl,dash pattern=on 8\pxl off 6\pxl]
            coordinates {(2,1)(14,1)};
          % Series named in place, upper left -- no legend box, no swatch of a
          % colour the curve does not already use.
          \node[anchor=west,font=\avlfAnnot\wmed,text=avlalert]
            at (axis cs:2.55,7.3) {\textcolor{avlalert}{\rule[0.6ex]{20\pxl}{2.4\pxl}}\ \ deletion --- $\lfloor h/2 \rfloor$};
          \node[anchor=west,font=\avlfAnnot\wmed,text=avlfib]
            at (axis cs:2.55,6.4) {\textcolor{avlfib}{\rule[0.6ex]{8\pxl}{2.4\pxl}\hskip4\pxl\rule[0.6ex]{8\pxl}{2.4\pxl}}\ \ insertion --- always 1};
        \end{axis}
      \end{tikzpicture}
    \end{column}
    \begin{column}{4.55cm}
      \eyebrow{One million nodes}\\[0.5em]
      {\color{avlrule}\hrule height 1\pxl}\vskip 8\pxl
      \avlstatrow[avlfib]{Insertion}{1}
      \avlstatrow[avlalert]{Deletion}{13}
      \vskip 0.6em
      {\avlfSec\color{avlmuted}Measured on a real AVL tree, every height
       from 2 to 14.}
    \end{column}
  \end{columns}

  \note{Insertion is flat at one, forever. Deletion grows with the height.
  At one million nodes the tree is 27 tall, so a single delete can cost 13
  rotations where an insert costs one.}
\end{frame}

\begin{frame}[label=p3]{Deletion: walk back to the root}
  \begin{center}
    {\avlfStmt\wbold\boldmath\color{avlink}BST deletion: $O(\log n)$.\par}

    \vspace{1.2em}
    {\avlfBody\color{avlmuted}After removing the node, walk upward\\
     from its parent toward the root.\par}

    \vspace{1.2em}
    \begin{avlpanelh}[0.78\textwidth]{avlpthree}
      \avlfLead\color{avlpthree}At most $h=O(\log n)$ ancestors are visited.\par
    \end{avlpanelh}

    \vspace{0.9em}
    {\avlfSec\color{avlmuted}Update the height. Check the balance. Move up.\par}
  \end{center}
  \note{First perform the usual BST deletion. The AVL height bound makes the
  search and successor or predecessor search logarithmic. Start retracing
  from the parent of the node physically removed, which may differ from the
  key initially requested when that key has two children. The path contains
  at most h ancestors, where h is the original tree height. With stored
  heights or balance factors, checking and updating each node costs O(1),
  apart from any rotations.}
\end{frame}

\begin{frame}[label=p3]{Deletion: at most two rotations per level}
  \begin{center}
    {\avlfBody\color{avlmuted}Each unbalanced node needs at most two rotations.\par}

    \vspace{1.0em}
    \begin{minipage}{0.88\textwidth}
      \begin{columns}[T,onlytextwidth]
        \begin{column}{0.47\textwidth}
          \begin{avlpanelh}[\linewidth]{avlfib}
            {\avlfLead\wbold\color{avlfib}LL / RR\par}
            \vspace{0.4em}
            {\avlfBody\color{avlbody}1 rotation\par}
          \end{avlpanelh}
        \end{column}
        \begin{column}{0.47\textwidth}
          \begin{avlpanelh}[\linewidth]{avlfib}
            {\avlfLead\wbold\color{avlfib}LR / RL\par}
            \vspace{0.4em}
            {\avlfBody\color{avlbody}2 rotations\par}
          \end{avlpanelh}
        \end{column}
      \end{columns}
    \end{minipage}

    \vspace{1.1em}
    {\avlfSec\color{avlmuted}At most $h$ ancestors, at most two rotations each:\par}
    \vspace{0.5em}
    {\avlfStmt\color{avlpthree}$\text{Max rotations}\le 2h=O(\log n)$\par}
  \end{center}
  \note{Count primitive rotations, so a double repair counts as two.
  LL and RR need one; LR and RL need two. Deletion may trigger repairs at
  multiple ancestors, but each visited level needs at most one single or
  double repair. Multiplying at most h ancestors by at most two rotations
  gives the safe upper bound 2h. It need not be tight, and many deletions
  require fewer rotations or none.}
\end{frame}

\begin{frame}[label=p3]{AVL deletion stays logarithmic}
  \begin{center}
    {\avlfBody\color{avlmuted}Each rotation takes $O(1)$. Add up the work:\par}

    \vspace{0.8em}
    {\avlfLead\color{avlbody}$\displaystyle
      \underbrace{O(\log n)}_{\text{BST deletion}}
      +\underbrace{O(h)}_{\text{upward walk}}
      +\underbrace{O(h)}_{\text{rotations}}$\par}

    \vspace{1.0em}
    {\avlfStmt\color{avlpthree}$\boxed{\text{AVL Deletion}=O(\log n)}$\par}

    \vspace{1.1em}
    \begin{avlpanelh}[0.96\textwidth]{avlpthree}
      \avlfSec\color{avlbody}Unlike insertion, deletion may rebalance at multiple levels.\\
      Even if repairs reach the root, the total stays $O(\log n)$.\par
    \end{avlpanelh}
  \end{center}
  \note{BST deletion costs O(log n). Visiting at most h ancestors costs O(h).
  At most 2h constant-time rotations also cost O(h), not O(h squared).
  Because h is O(log n), the sum is O(log n). The contrast with insertion is
  the number of repair sites: insertion stops after repairing one site,
  while deletion may continue repairing all the way to% c the root. That extra
  work does not change the logarithmic worst-case bound on the full operation.}
\end{frame}

\begin{frame}[label=p3]{Why bother with the worse constant}
  \begin{center}
  \begin{minipage}{0.90\textwidth}
    \avlfSec\color{avlmuted}
    {\fib{\wbold Databases and filesystems}}\par
    In-memory indexes need guaranteed worst-case lookup,
    not just a good average.\par\vspace{0.9em}
    {\fib{\wbold Read-heavy workloads}}\par
    Lookups vastly outnumber updates. The extra deletion work barely matters.\par\vspace{0.9em}
    {\fib{\wbold A hard bound beats a soft average}}\par
    Real-time systems care about the worst case, not the typical case.\par
  \end{minipage}

  \vspace{1.1em}
  {\avlfSec\color{avlbody}An unbalanced tree is fast until it isn't.
   An AVL tree is always fast\\--- that's the whole trade.\par}
  \end{center}
  \note{AVL trees are useful for ordered in-memory data, not a blanket choice
  for every database or filesystem index. OpenZFS uses AVL trees to track
  in-memory buffers; see include/sys/dnode.h, dn\_dbufs, in the OpenZFS source.
  Disk-oriented database indexes often use B-trees instead; see SQLite's
  Architecture of SQLite, B-Tree section. The read-heavy point is a workload
  trade-off, not a claim that AVL always beats other balanced trees.
  Always fast means logarithmically bounded tree work, not constant time or
  a guaranteed wall-clock deadline. Real-time systems also have to bound
  allocation, synchronization, and the cost of comparing keys.}
\end{frame}

% ------------------------------------------------------------------ closing
\avlstatement{A balanced tree is not one that is even.\\[0.35em]
It is one that is never more than \dstrong{one step} away from even.\\[0.35em]
That one step is what keeps every search fast.}{avlpthreeA}{avlpthreeB}{avlpthreeC}

\end{document}
